\documentclass[carta,firma,sinnotas]{lafrox}
\usepackage{fvextra}
\usepackage{pdflscape}

\datoslafrox{
  documento = {Subdocumento 8: Anexos},
  subtitulo = {Anexos 8.A--8.H},
  alcance = {Gestión de riesgos de la propuesta LafroX},
  formulario = {T-16: Plan de riesgos},
  fecha = {27 de Octubre de 2026},
  version = {0.1},
  nomenclatura = {LAFROX-Subdocumento8-Anexos},
}
% Soporte local de firma para el cuerpo y los anexos del Subdocumento 8.
% No altera la clase ni las plantillas compartidas.
\newcommand{\LFXportadasd}{%
  \AddToHookNext{shipout/foreground}{%
    \begin{tikzpicture}[remember picture,overlay]
      \node[anchor=north west,inner sep=0pt]
        at \LFXen{\LFXmi}{0.835}
        {\includegraphics[height=13mm]{logo/FIRMA_Alex_Aravena_transparente.png}};
    \end{tikzpicture}%
  }%
  \portadalafrox
}

% Firma abreviada sobre filete, en el mismo espacio del pie existente.
\renewcommand{\LFXpiederecho}{%
  \begingroup\LFXpiebase
  \begin{tabular}[b]{@{}r@{}}
    \raisebox{0pt}[0pt][0pt]{%
      \begin{tabular}[b]{@{}c@{}}
        \includegraphics[height=7mm]{logo/FIRMA_Alex_Aravena_transparente.png}\\[-1pt]
        \rule{3.4cm}{0.4pt}
      \end{tabular}%
    }\\[-1pt]
    {\intermedium\fontsize{9pt}{11pt}\selectfont\color{lafroxTinta}%
      Media firma: representante\quad\LFXfoliotexto}
  \end{tabular}%
  \endgroup
}


% Formato de portada común, basado en el Subdocumento 1.
\configurarportadalafrox{Anexos del Subdocumento 08}{Anexos del Subdocumento 08}{Anexos 8.A--8.H}{LAFROX-Subdocumento8-Anexos.pdf}{logo/FIRMA_Alex_Aravena_transparente.png}
\begin{document}
\portadalafrox
\preliminareslafrox
\setcounter{chapter}{7}

\chapter{Anexos del subdocumento 8}

Este documento reúne los anexos del plan de riesgos de LafroX.

\section{Anexo 8.A: Matriz de respuestas propuestas}

La Tabla 8.A.1 presenta una respuesta propuesta para cada uno de los 35 riesgos, con sus disparadores y responsables. La Tabla 8.A.2 especifica para los 28 riesgos con coordinación del CLIENTE la contraparte, solicitud, evidencia, plazo y contingencia por falta de respuesta. Sus objetivos de coordinación se proponen para acuerdo en el arranque; no acreditan una aceptación previa ni la eficacia de las respuestas. La estimación de esfuerzo por perfil, los plazos completos y la valoración residual proyectada están en 8.F.1--2; el costo-beneficio de las respuestas altas/críticas, en 8.F.3. La subcontratación no transfiere la obligación contractual final de LafroX.

\begingroup
\renewcommand{\thetable}{8.A.1}
\begin{landscape}
\setlength{\textwidth}{\linewidth}
\begin{tablalafrox}{Matriz de respuestas propuestas para los 35 riesgos vigentes. Fuente: riesgos, disparadores y responsables propuestos en el Formulario T-16; respuestas propuestas en este plan.}{tab:sd8-anexo-a-respuestas}{G{0.06} F{0.22} Y F{0.18}}{\cab{ID} & \cab{Disparador observable} & \cab{Respuesta propuesta y actuación ante el disparador} & \cab{Responsable propuesto en T-16}}
R-001 & Cadena del alcance sin certificación antes del límite de 28 días & Solicitar al Gerente Comercial, con TI y cada cadena, agenda EDI en mes 13, confirmada en mes 14, y certificación a D-5 del inicio de las cuatro semanas finales de evidencia de E2. Ante falta de respuesta, JP escala por CT; continuar pruebas independientes sin acreditar la cadena pendiente ni sustituir EDI por carga manual. & Arquitecto de Solución \\
R-002 & Lotes o vínculos históricos sin evidencia de origen al perfilar la carga & Solicitar a la Jefa de Calidad, con TI y Bodega, evidencia de origen y decisión sobre lotes no reconstruibles; política V-07 al cierre del mes 2 y excepciones a D-5 de validar cada carga. Si no responde, JP escala por CT, se aísla el conjunto y queda pendiente su aceptación; no se inventa historia sanitaria. & Líder de Datos \\
R-003 & V-09 no acredita tres años utilizables o revela períodos incompletos & Solicitar al Jefe de Abastecimiento y a TI el historial de tres años y validación de V-09 al cierre del mes 3; corregir extractos a D-5 de validar la carga. Ante silencio, escalar por CT, separar períodos incompletos y continuar el perfilado disponible sin reducir la historia exigida. & Líder de Datos \\
R-004 & Diferencias sin causa ni ajuste aprobado en el acta previa al corte & Solicitar a Bodega, con Operaciones y Finanzas, protocolo de conciliación en mes 2, validado en mes 3. Confirmar conteo a D-5 del corte y cerrar acta y ajustes antes de D. Si no hay respuesta, escalar por CT, repetir conciliación y bloquear el corte afectado hasta contar con evidencia aprobada. & Líder de Datos \\
R-005 & No existe saldo inicial aceptado por cliente y transportista & Solicitar al Gerente de Operaciones y a Finanzas respaldos y conciliación V-08 de saldos por cliente y transportista; política en mes 2 y respuesta a D-5 de validar la carga M8 de cada ola. Ante silencio, escalar por CT y mantener M8 pendiente para esos saldos; no presumir saldo cero. & Líder de Datos \\
R-006 & V-10 no obtiene horas o montos históricos verificables & Solicitar a Comercial y Operaciones las horas históricas y a Finanzas los montos de rendición, con validación de V-10 al cierre del mes 1. Si no responden, escalar por CT y mantener pendiente la referencia de aceptación; la captura futura no reconstruye la historia anterior. & Líder de Implantación y Gestión del Cambio \\
R-007 & Documentos todavía vigentes quedan fuera del volumen histórico previsto & Solicitar a Finanzas y TI el universo de deuda vigente, con respuesta inicial al cierre del mes 3 y conciliación a D-5 de validar la migración. Ante silencio, escalar por CT y mantener la cartera incompleta pendiente; preparar extracción complementaria sin omitir deuda obligatoria. & Líder de Datos \\
R-008 & Trámite tributario permanece pendiente de respuesta externa & Mantener trazabilidad de trámites pendientes, consulta y reintentos idempotentes hasta recibir confirmación del SII. Ante demora externa, conservar el estado pendiente y escalar según el procedimiento acordado; el silencio no equivale a aprobación. & Líder de Operación / SRE \\
R-009 & Calidad no aprueba parámetros por producto al preparar certificación & Solicitar a la Jefa de Calidad la matriz térmica versionada V-11 durante mes 1 y su aprobación al cierre del mes 2, antes de certificar. Si no responde, escalar por CT y mantener el tratamiento preventivo de lecturas fuera de ficha; no cerrar marcha blanca ni aprobar parámetros por el CLIENTE. & Líder de Calidad \\
R-010 & Equipo de terreno especificado no recibido antes de instalar o probar & Solicitar al Jefe de TI, con Operaciones, calendario de equipos en mes 1, entrega confirmada a D-5 y recepción técnica antes de la instalación o prueba D. Ante silencio o faltante, escalar por CT, continuar trabajo independiente y bloquear las actividades sin equipo. & Líder de Implantación y Gestión del Cambio \\
R-011 & Acuerdo operativo de un transportista pendiente antes de su ola & Solicitar al Gerente de Operaciones y al representante de cada transportista el acuerdo y la nómina desde mes 2, firmados a D-5 de su incorporación a la ola. Ante silencio, escalar por CT y no habilitar a esa empresa; conservar su cobertura pendiente sin retirarla del alcance. & Líder de Implantación y Gestión del Cambio \\
R-012 & Propietario niega acceso o autorización de instalación del termógrafo & Solicitar a Operaciones y al propietario permiso de instalación por vehículo con D-10 de la visita y obtenerlo a D-5. Ante silencio, escalar por CT, suspender la intervención y proponer otra visita o vehículo autorizado; mantener pendiente la evidencia térmica. & Líder de Implantación y Gestión del Cambio \\
R-013 & Transacciones se siguen registrando solo por el procedimiento anterior & Solicitar a Comercial y Operaciones, con Bodega por turno, responsables y agenda de acompañamiento a D-5 del inicio de cada ola, convocando a D-10. Si no responden o persiste el registro externo, escalar por CT, focalizar acompañamiento y mantener pendiente la evidencia de adopción. & Líder de Implantación y Gestión del Cambio \\
R-014 & Planificador deja de estar disponible antes de validar reglas pendientes & Solicitar al Gerente de Operaciones y al planificador confirmación de retiro y disponibilidad V-14 al cierre del mes 1; transferir y validar reglas con el ayudante en meses 1--3. Probar las rutas y excepciones sin el planificador en la marcha blanca de meses 13--15 (S7, 7.2.1). Si falta respuesta o se retira antes, escalar por CT y reconstruir con ayudante e historia verificable, sin dar por validadas reglas pendientes. & Líder de Implantación y Gestión del Cambio \\
R-015 & Usuarios convocados no completan formación y evaluación previstas & Solicitar a Operaciones, Bodega, Comercial y TI, según perfil, nómina y liberación horaria a D-5 de cada sesión convocada a D-10. Ante silencio o ausencia, escalar por CT y reprogramar evaluación; usar sólo suplentes certificados y no acreditar preparación ni cierre de ola sin evidencia. & Líder de Implantación y Gestión del Cambio \\
R-016 & Decisión requerida no llega en la fecha acordada y bloquea la tarea & Solicitar a CT área decisora y suplente, con agenda al cierre del mes 1; cursar insumos a D-10 y pedir respuesta a D-5. Las aceptaciones formales conservan diez días hábiles. Para promover la región de contingencia, avisar al suplente a los 15 min sin respuesta del titular y obtener decisión dentro de 30 min (S4, 4.3.2.5); no esperar el plazo ordinario. Mantener bloqueada la acción sin autorización; el silencio no permite cambios. & Jefe de Proyecto \\
R-017 & Corte del enlace del sitio o del respaldo restante; en Los Ángeles de madrugada, pérdida del camino Starlink disponible. & Mitigar. Verificar respaldos antes de aceptar la ola de Concepción; no considerar disponible un enlace sólo por estar ofertado. Mantener y ensayar los enlaces ya ofertados: en Talca y Concepción, fibra y LTE con Starlink en espera caliente; en cross-docking, Starlink con LTE Entel/Claro y conmutación menor a 30 s. En Los Ángeles de madrugada, probar pérdida del camino Starlink restante y continuidad local de su ventana de 3 h. Coordinar restitución y conciliación sin duplicar el mismo corte como falla de SII, POS o mapas. Referencia: S4, 4.2.5.3. & Líder de Operación / SRE \\
R-018 & Corte eléctrico excede autonomía del respaldo aplicable & Verificar autonomía, alarmas y procedimiento de recuperación de los componentes de LafroX, coordinando las dependencias de energía con el CLIENTE/proveedor. Si se excede la autonomía, activar recuperación y registrar la indisponibilidad e impacto. & Líder de Operación / SRE \\
R-019 & Cliente no otorga aprobación de tratamiento en las regiones propuestas & Solicitar por CT, con TI y la autoridad del CLIENTE encauzada por Gerencia General, aprobación escrita de residencia y transferencia en mes 1, resuelta al cierre del mes 2 o antes de habilitar datos, lo que ocurra primero. Ante silencio, escalar y preparar alternativa de arquitectura; mantener deshabilitada la transferencia no aprobada. & Arquitecto de Solución \\
R-020 & Intervención urgente requiere autorización y esta no se obtiene & Solicitar a CT, Operaciones y TI flujo y suplentes al cierre del mes 1; ante incidente, pedir autorización escrita antes del cambio. Sin respuesta, escalar en una hora o antes si lo exige el plazo de atención, manteniendo contención autorizada y estado seguro. El silencio no permite intervenir durante congelamiento. & Encargado de Seguridad de la Información \\
R-021 & Falla mecánica impide completar el recorrido de prueba & Detener de forma segura el recorrido afectado, registrar los casos de prueba incompletos y coordinar con el CLIENTE/transportista una nueva ventana antes de repetir la validación. & Líder de Implantación y Gestión del Cambio \\
R-022 & Siniestro impide continuar los casos de prueba previstos & Priorizar la atención del accidente conforme a los procedimientos del transportista y del CLIENTE; suspender la prueba, registrar su impacto y reprogramar sólo cuando exista autorización para continuar. & Líder de Implantación y Gestión del Cambio \\
R-023 & Asalto interrumpe recorrido y evidencia prevista & Coordinar con el CLIENTE y el transportista la respuesta al asalto, robo o incendio, proteger a las personas y registrar qué evidencia y pruebas se perdieron. Acordar una nueva ejecución segura antes de continuar la validación. & Líder de Implantación y Gestión del Cambio \\
R-024 & Conductor sustituto no figura autorizado/certificado al presentarse & Solicitar a Operaciones y al transportista nómina y suplentes a D-5 de incorporar la empresa y confirmar conductor, vehículo y viaje antes de cada salida. Si falta respuesta o llega un sustituto, escalar, verificar en andén y solicitar conductor certificado; no habilitar uso no autorizado ni acreditar cobertura incompleta. & Líder de Implantación y Gestión del Cambio \\
R-025 & Solicitud nueva excede requisitos o exclusiones aprobados & Solicitar por CT a la gerencia peticionaria solicitud escrita y decisión del cambio antes de su primera actividad dependiente; registrar y remitir el mismo día hábil. Ante silencio, mantener línea base y continuar trabajo independiente; si amenaza un hito, escalar al Comité Ejecutivo sin iniciar desarrollo no aprobado. & Jefe de Proyecto \\
R-026 & Cambio aplicable obliga a evaluar adecuación de controles o retención & Revisar cambios normativos aplicables, determinar su efecto con las áreas competentes y tramitar ajustes de controles, retención o eliminación mediante decisión y control de cambios documentados, sin posponer una obligación legal aplicable. & Arquitecto de Solución \\
R-027 & Diferencia de cálculo o reclamo en una liquidación mensual & Solicitar a Finanzas, con Comercial y CT, reglas de 5.4.3 aprobadas al cierre del mes 20 y tres liquidaciones en sombra reproducidas en meses 21--23 antes del pago variable del mes 24. Ante silencio, no activar el tramo; una disputa se escala y suspende sólo el devengo correspondiente según contrato, manteniendo el servicio. & Jefe de Proyecto \\
R-028 & Acuerdo de finalidades y controles pendiente a D-5 de la ola, u objeción formal al uso operacional propuesto. & Mitigar y escalar. JP coordina con Operaciones, TI y sindicato el acuerdo 5.4.2 antes de la ola: convocar a D-10 y resolver a D-5. Explicar finalidades, acceso y controles; mantener excluidas las cámaras en cabina y no usar GPS para jornada sin acuerdo previo. Ante desacuerdo, escalar por CT, mantener sólo el uso operacional autorizado y registrar la cobertura de la ola pendiente. & Jefe de Proyecto \\
R-029 & Preparador nuevo asignado sin tutor o evaluación, o aprendizaje que excede dos horas en el programa 7.1.2. & Mitigar. Mantener 7.1.2 en meses 13--56, turno 22:00--06:00, tutor y evaluación antes del trabajo autónomo, con aprendizaje máximo de dos horas. Bodega y Operaciones entregan altas por turno antes de asignarlas. IMP conduce hasta mes 21 y SRE desde 22, familia IMP. Ante sobrecarga, usar personal certificado, reforzar la inducción y escalar por CT; no duplicar las HH ordinarias de 7.1.2. & Líder de Implantación y Gestión del Cambio \\
R-030 & No hay contrato de interfaz o extracción verificado al cierre del mes 1, o el ensayo temprano contradice formato, acceso o comportamiento declarados. & Mitigar y escalar. ARQ solicita a TI acceso, diccionario y condiciones de extracción en mes 1 (1.2.3, 3.3.2 y V-16); DAT verifica la migración del WMS. Probar temprano sin modificar el ERP. Si falta una interfaz, evaluar intercambio de archivos controlado, sujeto a demostrar los requisitos; elevar por CT el bloqueo y sus impactos. La alternativa no acredita por sí sola latencia, integridad ni emisión tributaria. & Arquitecto de Solución \\
R-031 & Incidente alto/crítico, evidencia o certificación pendiente en las semanas finales del peak, sin ventana de intervención o validación disponible. & Mitigar y escalar. JP acuerda con Operaciones y CT disponibilidad y continuidad antes de septiembre; cerrar habilitaciones en agosto según S7. En meses 19--20 mantener evidencia diaria de las seis condiciones de marcha blanca. Un defecto durante 1--25 de septiembre no autoriza despliegue: activar continuidad aprobada y R-020 si requiere autorización. Si persiste una condición incumplida, no declarar cierre; elevar plan de recuperación sin mover hitos por cuenta propia. & Jefe de Proyecto \\
R-032 & Aviso de soporte/compatibilidad que deja una versión necesaria fuera de la política aprobada o hace fallar la regresión de la actualización prevista. & Mitigar. ARQ mantiene inventario de versiones y revisa soporte y compatibilidad trimestralmente desde mes 1, con DES y SRE. Preparar y probar adecuación antes del límite verificado de soporte, dentro de ventanas autorizadas. Si el cambio excede la actualización ordinaria 8.2.5, registrar el esfuerzo extraordinario y elevar impactos por CT; usar reversión ensayada y soporte permitido, sin sostener una versión insegura por conveniencia. & Arquitecto de Solución \\
R-033 & Ensayo de exportación/restauración incompleto, acceso o clave no disponible para el CLIENTE, o cambio de condiciones que impide cumplir el plan de salida. & Mitigar. ARQ verifica con TI propiedad de cuentas, acceso a datos, claves, exportaciones e infraestructura como código en la recepción inicial y en cada actualización anual del plan 8.1.4. Probar una restauración utilizable, no sólo descargar archivos. Ante restricción, documentar brecha, mantener continuidad autorizada y escalar por CT alternativas y decisión contractual; no declarar reversibilidad completa ni cambiar de proveedor sin prueba y aprobación. & Arquitecto de Solución \\
R-034 & Carga sin guía vigente que requiere nueva emisión mientras el ERP está indisponible; pérdida de sala con ERP aún no restituido. & Mitigar y escalar. ARQ/SRE acuerdan con TI y Finanzas restitución del ERP y ensayan AL-DTE-01/AL-DR-01 antes de H5/H10, incluida pérdida de sala. Mantener sólo salidas con guía preemitida aún válida y bloquear las que requieren nueva guía. Recuperar WMS en nube no restaura el ERP ni exceptúa el RTO requerido; registrar y escalar por CT cualquier incumplimiento, sin segundo emisor ni aprobación tributaria por silencio. & Arquitecto de Solución \\
R-035 & Conductor asignado y habilitado rechaza el flujo digital o entrega sin la evidencia requerida, sin que la causa sea firma pendiente o reemplazo sin aviso. & Mitigar y escalar. IMP coordina con Operaciones, transportista y conductor demostración y acompañamiento antes de su salida, dentro de 7.1.3. Confirmar disponibilidad a D-5 de la ola y aceptación práctica en andén. Ante rechazo, solicitar a la empresa conductor certificado y dispuesto, mantener acompañamiento y elevar por CT la cobertura pendiente; no atribuir a LafroX facultades laborales ni declarar completa la evidencia faltante. & Líder de Implantación y Gestión del Cambio \\
\end{tablalafrox}
\end{landscape}
\endgroup

\subsection*{Dependencias del CLIENTE: compromisos y contingencias}
\addcontentsline{toc}{subsection}{Dependencias del CLIENTE: compromisos y contingencias}

Esta matriz establece los compromisos de coordinación de los 28 riesgos que requieren insumos o decisiones del CLIENTE. La contraparte aporta antecedentes y decide dentro de su autoridad; LafroX conserva la responsabilidad de diseño, capacidad, seguridad y mantenimiento. La gestión de servicios externos, suministros, siniestros y cambios normativos se realiza mediante las respuestas específicas de la Tabla 8.A.1. Los plazos de coordinación se proponen para acuerdo en el arranque; no acreditan respuestas recibidas ni modifican las obligaciones contractuales.

CT es la Contraparte Técnica. En ``ID / dueño'', JP, ARQ, DAT, CAL, IMP y SEG corresponden a las familias y responsables de la Tabla 8.1; desde mes 22 se aplica el relevo allí definido. D identifica la actividad o ventana indicada de la línea base; D-10 y D-5 son días hábiles anteriores. Salvo indicación específica, el dueño envía expediente completo a D-10 y solicita respuesta a D-5; si existe un plazo de mes más temprano, rige ese plazo. Para mes 1 solicita en sus primeros cinco días hábiles. Los plazos formales del art. 18.3 se conservan.

Si falta respuesta, el dueño informa a JP y CT el mismo día hábil del vencimiento. JP pide encauzamiento en un día hábil y, si no se resuelve en dos días hábiles del vencimiento, solicita instancia ejecutiva extraordinaria con la gerencia competente. Las amenazas a seguridad, continuidad o hitos se escalan de inmediato. La columna final añade la contingencia específica; no se espera la reunión mensual para actuar. Toda respuesta se cierra con evidencia validada, no con una promesa o el envío de la solicitud. Los titulares y suplentes por cargo se registran en mes 1 y se actualizan al cambiar su designación.

\begingroup
\renewcommand{\thetable}{8.A.2}
\begin{landscape}
\setlength{\textwidth}{\linewidth}
\begin{tablalafrox}{Contrapartes, solicitudes, plazos y contingencias de los 28 riesgos con coordinación del CLIENTE. Fuente: asignaciones de S2/S3/S6/S13 y objetivos de coordinación propuestos en S8, 8.1.2.}{tab:sd8-dependencias-cliente}{G{0.07} F{0.17} Y F{0.21} F{0.23}}{\cab{ID / dueño} & \cab{Contraparte del CLIENTE / apoyo} & \cab{Solicitud y evidencia de resolución} & \cab{Plazo de solicitud y respuesta} & \cab{Escalamiento y contingencia por falta de respuesta}}
R-001 / ARQ & Gerente Comercial; coordinación de TI con el responsable EDI de cada cadena. & Confirmar cadenas del alcance, contacto técnico, agenda y ambiente de certificación; entregar resultado firmado de intercambio de pedidos y avisos por cadena. & Solicitar agenda en mes 13 y confirmarla al cierre del mes 14. Certificación a D-5 de la ventana final de 28 días de evidencia de E2; D es el inicio de esa ventana. & Escalar al Gerente Comercial por CT. Reordenar pruebas independientes; conservar pendiente la cobertura EDI de la cadena. La carga manual no acredita E2 ni permite excluirla. \\
R-002 / DAT & Jefa de Calidad; Jefe de TI y Jefa de Bodega aportan fuentes ERP/WMS y antecedentes de lotes. & Entregar extractos y evidencia de origen de los vínculos sanitarios; Calidad debe decidir por escrito la disposición de conjuntos cuya historia no pueda reconstruirse y la política ante lote desconocido (V-07). & Solicitar fuentes en mes 1; política y resolución inicial al cierre del mes 2. Para cada carga, resolver las excepciones a D-5 de su validación; nunca después del corte. & Escalar a CT con los lotes y evidencia faltante; aislar el conjunto y trabajar sobre datos verificables. Mantener la aceptación pendiente; no inventar lotes ni reconstruir historia sin fuente. \\
R-003 / DAT & Jefe de Abastecimiento; Jefe de TI entrega y explica los extractos del ERP. & Entregar tres años de ventas y pedidos, diccionario y cobertura temporal; validar el informe de perfilado V-09 y las excepciones identificadas. & Solicitar en mes 2; validación de V-09 al cierre del mes 3 y correcciones de cada extracto a D-5 de validar su carga. & Escalar por CT a Abastecimiento y TI; separar períodos incompletos y continuar perfilado de fuentes disponibles. No reducir unilateralmente los tres años exigidos ni dar por aceptado el historial faltante. \\
R-004 / DAT & Jefa de Bodega; Gerente de Finanzas aprueba ajustes de saldo y Gerente de Operaciones coordina conteos. & Entregar conciliación entre inventario legado y conteo físico, causas de diferencias y ajustes autorizados; presentar acta de validación previa al corte. & Solicitar protocolo en mes 2 y validarlo en mes 3. En cada corte, responsables y conteo confirmados a D-5; conciliación y ajustes firmados antes de D. & Escalar por CT a Bodega y Finanzas; repetir conteo y conciliación del conjunto afectado. No ejecutar su corte sin acta; continuar trabajo independiente y evaluar recuperación sin mover hitos por silencio. \\
R-005 / DAT & Gerente de Operaciones; Gerente de Finanzas valida saldos y representantes de transportistas aportan antecedentes. & Entregar saldos de envases por cliente y transportista y sus respaldos; validar V-08. Si se requiere conteo inicial, autorizar su ejecución y firmar conciliación de M8. & Solicitar política de saldos en mes 2 (S-10). Para cada ola de reparto: respuesta a D-5 de validar la carga M8 y conciliación aprobada antes de habilitarla. & Escalar a Operaciones y Finanzas por CT; preparar conteo o reconstrucción documental autorizados. Mantener pendiente M8 en los saldos afectados; no asignar cero ni presumir aceptación. \\
R-006 / IMP & Gerentes Comercial y de Operaciones: horas de entrega; Gerente de Finanzas: montos esperados de rendición. & Entregar registros históricos verificables y aprobar fuente, período y confiabilidad de la línea base V-10; identificar limitaciones y resolver formalmente una alternativa que preserve los criterios de aceptación. & Solicitar en los primeros cinco días hábiles del mes 1; validar V-10 al cierre de ese mes. Cualquier alternativa debe resolverse antes de iniciar la comparación de aceptación. & Escalar a las tres gerencias por CT. Conservar por separado la captura futura y la historia comprobada; suspender la comparación afectada hasta resolver la referencia, sin fabricar antecedentes ni cambiar la meta. \\
R-007 / DAT & Gerente de Finanzas; Jefe de TI ejecuta la extracción del ERP y documenta su cobertura. & Entregar listado completo de deuda vigente y conciliación con el extracto provisional; autorizar la inclusión de documentos omitidos y registrar volumen y diferencias. & Solicitar universo y criterio de vigencia en mes 2; respuesta inicial al cierre del mes 3. Cada extracto completo y conciliado a D-5 de validar la migración de cartera. & Escalar por CT a Finanzas y TI; medir y preparar la extracción complementaria. Mantener pendiente la cartera incompleta; tramitar impactos por el gobierno del proyecto sin tratar la integridad obligatoria como ampliación opcional. \\
R-009 / CAL & Jefa de Calidad, autoridad de aprobación sanitaria; Bodega y Operaciones validan aplicación en terreno. & Aprobar y versionar matriz de umbrales y duraciones por tipo de producto, fichas de almacenamiento y reglas de liberar, bloquear o rechazar; emitir aprobación de V-11. & Solicitar durante mes 1; matriz aprobada al cierre del mes 2, antes de certificar parametrización. Todo cambio posterior debe aprobarse antes de probarlo o usarlo. & Escalar inmediatamente por CT si amenaza certificación. Aplicar el tratamiento preventivo de V-11 a lecturas fuera de ficha y conservar el bloqueo; no cerrar marcha blanca sin matriz aprobada. \\
R-010 / IMP & Jefe de TI, como coordinador técnico de adquisición; Gerente de Operaciones confirma disponibilidad por sitio y ola. & Confirmar orden y fecha de entrega del equipamiento comprado por el CLIENTE, inventario por ola y evidencia de recepción técnica conforme a la especificación de LafroX. & Solicitar calendario en mes 1. En cada instalación, entrega y acceso confirmados a D-5 y recepción técnica antes de D; D es el inicio de instalación o prueba que requiere el equipo. & Escalar por CT al Jefe de TI y a Operaciones; ejecutar sólo trabajo independiente. Bloquear instalación y pruebas sin equipo; evaluar recuperación sin sustituir la compra del CLIENTE por una obligación de LafroX. \\
R-011 / IMP & Gerente de Operaciones; representante autorizado de cada empresa transportista firma su acuerdo. & Obtener acuerdos operacionales por empresa con nómina y coordinación de conductores, vehículos, acceso y capacitación; registrar cobertura de las diez empresas sin presumir firma. & Solicitar acuerdos en mes 2. Para cada ola de reparto, firma y nómina a D-5 del inicio de incorporación de la empresa; verificar antes de habilitar a sus conductores. & Escalar por CT a Operaciones; incorporar sólo empresas con acuerdo y registrar cobertura pendiente. Una ola sin esa empresa no acredita cobertura completa ni elimina obligaciones de aceptación. \\
R-012 / IMP & Gerente de Operaciones gestiona acceso; propietario o representante autorizado del camión otorga permiso. & Entregar autorización de acceso e instalación de termógrafo por vehículo, agenda de visita y condiciones del propietario; documentar cualquier alternativa sometida a aprobación. & Solicitar con D-10 y recibir autorización a D-5 de cada visita de instalación. Comprobar permiso y vehículo antes de intervenir en D. & Escalar a Operaciones por CT; cancelar la intervención no autorizada y proponer otra visita o vehículo autorizado con trazabilidad equivalente. Mantener pendiente la evidencia térmica; no instalar ni sustituir controles sin aprobación. \\
R-013 / IMP & Gerentes Comercial y de Operaciones; Jefa de Bodega para el turno nocturno y supervisores de cada ola. & Confirmar supervisores, turnos y usuarios que registrarán transacciones oficiales; aprobar agenda de acompañamiento y revisar evidencia diaria de adopción por usuario y ola. & Convocar con D-10 y confirmar responsables y agenda a D-5 del inicio de cada ola. Revisar transacciones diariamente durante marcha blanca y resolver ausencias el mismo día. & Escalar por CT a la gerencia del turno o canal; focalizar acompañamiento y registrar brechas. No acreditar adopción con registros sólo en el proceso anterior ni cerrar la ola sin evidencia. \\
R-014 / IMP & Gerente de Operaciones; Planificador de Rutas y su ayudante participan en transferencia y validación. & Confirmar fecha efectiva de retiro y disponibilidad (V-14), reservar talleres y entregar reglas y excepciones documentadas con prueba de transferencia al ayudante. La transferencia se comprueba con rutas y excepciones probadas sin el planificador durante la marcha blanca. & Solicitar disponibilidad en los primeros cinco días hábiles del mes 1 y confirmarla al cierre del mes 1. Ejecutar y validar transferencia en meses 1--3, antes de perder disponibilidad. Prueba sin el planificador en meses 13--15 (S7, 7.2.1). & Escalar por CT a Operaciones; activar talleres con ayudante y antecedentes verificables. Mantener reglas no corroboradas como incertidumbre y obtener validación antes de usarlas en rutas aceptadas. \\
R-015 / IMP & Gerente de Operaciones y Jefa de Bodega para turnos; Gerente Comercial para preventistas; Jefe de TI para administradores. & Confirmar nómina por perfil, liberación horaria, suplentes y acceso a sesiones; entregar disponibilidad para formación y evaluación y revisar evidencia individual de certificación. & Convocar con D-10 de cada sesión y confirmar nómina y horarios a D-5. Completar evaluación antes de habilitar al usuario y, en todo caso, antes de cerrar la marcha blanca de su ola. & Escalar por CT a la jefatura responsable el mismo día del incumplimiento; reprogramar sesiones y disponer suplente sólo si está certificado. Mantener pendiente la certificación y el cierre de ola afectado. \\
R-016 / JP & Contraparte Técnica, responsable de encauzar la decisión hacia Calidad, TI, Operaciones, Comercial o Finanzas según la materia. & Identificar en cada solicitud el área decisora, titular funcional y suplente, pregunta, antecedentes completos, tarea bloqueada y fecha límite; obtener respuesta o acta de decisión de la autoridad competente. Para promoción regional, identificar autorizador titular y suplente con facultad delegada. & Registrar agenda y responsables al cierre del mes 1. Para insumos de coordinación, cursar a D-10 y recibir respuesta a D-5 de la tarea D; para recepción de entregables rigen los diez días hábiles del art. 18.3. Excepción de S4, 4.3.2.5: avisar al suplente a los 15 min sin respuesta del titular y resolver dentro del máximo de 30 min. & JP notifica a CT el vencimiento el mismo día, solicita encauzamiento en un día hábil y, si sigue sin resolución en dos días hábiles, pide instancia ejecutiva extraordinaria. En promoción regional, escalar al suplente a los 15 min y mantener sin autorizar la promoción si vence el máximo de 30 min; registrar el efecto sobre el RTO, sin promover por silencio. Si amenaza seguridad, continuidad o hito, escala inmediatamente y mantiene bloqueada la tarea dependiente. \\
R-019 / ARQ & Contraparte Técnica; Jefe de TI coordina revisión técnica y la Gerenta General encauza autorización con la asesoría competente del CLIENTE. & Resolver por escrito regiones, residencia, transferencia y tratamiento de datos sobre el expediente de arquitectura y seguridad; identificar autoridad firmante y condiciones de habilitación. & Solicitar expediente y autoridad en mes 1, decisión inicial al cierre del mes 2 y siempre antes de habilitar entornos con datos sujetos a autorización; aplicar la fecha que venza primero. & Escalar por CT a la autoridad identificada y al Comité de Arquitectura; preparar alternativa con sus impactos. Mantener deshabilitada la transferencia no aprobada; no elegir unilateralmente otra región ni presumir autorización. \\
R-020 / SEG & Contraparte Técnica; Gerente de Operaciones autoriza la intervención operacional y Jefe de TI valida su ejecución técnica. & Aprobar flujo de autorización urgente, contactos y suplentes; por incidente, decidir por escrito la intervención, ventana, alcance, reversión y responsable con el expediente de seguridad. & Flujo aprobado al cierre del mes 1. Cada intervención se solicita al detectar la necesidad; si no hay respuesta, reiterar y escalar en una hora o antes si el plazo de atención o continuidad exige actuar. Autorización siempre previa al cambio. & Aplicar contención permitida por el procedimiento, notificar inmediatamente a CT y Operaciones y convocar decisión ejecutiva urgente. Mantener estado seguro sin eludir congelamiento, alterar plazos contractuales ni interpretar silencio como permiso. \\
R-024 / IMP & Gerente de Operaciones y representante de la empresa transportista; supervisor de andén verifica al conductor real. & Confirmar por despacho conductor, vehículo y viaje; acreditar identidad, autorización y certificación del reemplazo y registrar la vinculación auditada persona--vehículo--viaje. & Pedir nómina y suplentes a D-5 de incorporar la empresa. En cada despacho, confirmar conductor antes de salida; reemplazo no avisado se verifica y vincula en andén antes de habilitar uso. & Escalar de inmediato al supervisor de Operaciones y transportista por CT; solicitar conductor certificado y realizar formación extraordinaria del sustituto. No habilitar acceso no autorizado ni contabilizar pruebas incompletas como aceptadas. \\
R-025 / JP & Contraparte Técnica; gerencia solicitante fundamenta la petición y Comité Ejecutivo decide los cambios dentro de sus atribuciones. & Entregar solicitud escrita con objetivo y criterio de aceptación; decidir alcance, esfuerzo y programación sobre el análisis de impacto antes de incorporar el cambio a la línea base. & Registrar y remitir a CT el mismo día hábil de recibirla; comunicar fecha de decisión necesaria antes de su primera actividad dependiente. No iniciar la implementación hasta aprobación formal. & Escalar al Comité Ejecutivo si la decisión amenaza un hito. Mantener la línea base vigente y continuar el trabajo independiente; el silencio no aprueba el cambio ni permite cargar desarrollo adicional. \\
R-027 / JP & Gerente de Finanzas, con Gerente Comercial para canasta y malla; Contraparte Técnica canaliza conformidad y controversias. & Aprobar canasta, índices y regla estacional de 5.4.3, reproducir tres liquidaciones en sombra y conciliar evidencia de cada liquidación y de las diferencias del período. & Acordar reglas en meses 17--20, antes de cerrar mes 20. Validar tres liquidaciones en sombra en meses 21--23 antes del primer efecto de pago del mes 24; revisión posterior antes de cerrar cada liquidación. & Escalar por CT a Finanzas y al Comité Ejecutivo. Sin regla o reproducción válida no activar el tramo; ante disputa, suspender sólo el devengo del período según el mecanismo contractual y mantener servicio y acompañamiento fijo. \\
R-028 / JP & Gerente de Operaciones y Jefe de TI; sindicato de conductores como interlocutor del acuerdo, por CT. & Acta sobre finalidades, acceso y controles de la telemetría y uso de terminales; documentar observaciones y el acuerdo exigible para cualquier control de jornada. & Convocar a D-10 y resolver a D-5 del inicio de la ola de reparto; acuerdo 5.4.2 antes de habilitar el uso que lo requiere. & JP escala por CT a Operaciones; no habilitar cámaras en cabina ni control de jornada sin acuerdo. Mantener el uso operacional autorizado y la cobertura pendiente identificada. \\
R-029 / IMP & Jefa de Bodega y Gerente de Operaciones; comunicación de altas y disponibilidad de tutores por turno. & Nómina de altas, turnos y tutores; registro de tiempo de aprendizaje y evaluación individual del preparador, con cobertura certificada mientras completa formación. & Meses 13--56; informar cada alta antes de asignarla al turno y completar inducción/evaluación antes del trabajo autónomo. Revisar cobertura en cada relevo. & IMP, y SRE desde mes 22, solicita cobertura certificada y refuerzo por CT a Bodega/Operaciones; mantener sin habilitación autónoma al preparador no evaluado. \\
R-030 / ARQ & Jefe de TI; custodios del ERP y WMS y proveedor del legado cuando su acceso sea necesario. & Accesos autorizados, formatos, diccionario y contratos de interfaz/extracción; resultado de ensayo temprano y restricciones del proveedor del ERP/WMS. & Solicitar en los primeros cinco días hábiles del mes 1 y verificar contrato inicial al cierre del mes 1; probar cada interfaz antes de su integración y cada extracción antes de validar la carga. & ARQ escala por CT a TI, prepara alternativa de archivo controlado y mide sus límites; mantener bloqueada la aceptación dependiente sin modificar el ERP ni omitir datos. \\
R-031 / JP & Gerente de Operaciones y Contraparte Técnica; gerencias usuarias para disponibilidad de validación. & Agenda de usuarios, decisiones y continuidad autorizada para el peak; estado diario de las condiciones de cierre y resolución formal de impactos en hitos. & Preparar antes del mes 19 y completar habilitaciones en agosto bajo el supuesto de inicio de S7. Verificación diaria en meses 19--20 y antes del cierre previsto del mes 20. & JP escala inmediatamente por CT al Comité Ejecutivo; mantener continuidad autorizada y no cerrar con condiciones incumplidas ni desplegar en congelamiento sin autorización. \\
R-032 / ARQ & Jefe de TI y Contraparte Técnica; Operaciones para ventanas de actualización y aceptación del servicio. & Calendario de ventanas, política de versiones y conformidad con las pruebas de adecuación; aprobación de cambios que afecten compromisos. & Inventario en mes 1, revisión trimestral hasta mes 56; acordar intervención antes del límite de soporte comprobado y antes de activar la versión. & ARQ escala por CT el impacto; ensayar reversión y reprogramación autorizada de la adecuación, preservando atención de vulnerabilidades y compromisos de servicio. \\
R-033 / ARQ & Jefe de TI y Contraparte Técnica; autoridad contractual del CLIENTE para condiciones de salida y acceso. & Accesos y titularidad de cuentas/claves, evidencia de exportación y restauración y decisión sobre restricciones contractuales de sustitución. & Verificar en recepción inicial y actualización anual de 8.1.4; resolver cada restricción antes del ensayo o transición de salida que dependa de ella. & ARQ eleva por CT la restricción al Comité Ejecutivo y prepara salida gradual comprobable; conservar acceso y continuidad bajo condiciones autorizadas sin presumir equivalencia entre proveedores. \\
R-034 / ARQ & Jefe de TI y Gerente de Finanzas; Operaciones coordina cargas y salidas afectadas. & Procedimiento de restitución del ERP, responsables y evidencias de emisión/anulación, con resultado de ensayo de pérdida de sala y cambio de carga. & Acordar procedimiento antes del ensayo AL-DTE-01/AL-DR-01 de H5/H10. En incidente, comunicar de inmediato a TI/Finanzas y decidir cada salida antes de despachar. & ARQ/SRE escalan por CT a TI/Finanzas; bloquear salida sin guía válida y conservar la carga original sólo si conserva guía vigente. Medir restitución completa frente al RTO exigido. \\
R-035 / IMP & Gerente de Operaciones y representante de cada empresa transportista; conductor como participante de la validación práctica. & Disponibilidad de conductores y demostración de uso en andén, evidencia de inducción y compromiso operativo del transportista ante rechazo individual. & Confirmar participantes a D-5 de la ola de reparto y demostrar uso antes de la primera salida del conductor; actuar en andén ante cada rechazo. & IMP escala por CT a Operaciones y al transportista; solicitar conductor certificado y dispuesto y acompañar. Registrar entregas/evidencias pendientes sin dar cobertura por cumplida. \\
\end{tablalafrox}
\end{landscape}
\endgroup

\clearpage
\section{Anexo 8.B: Contingencia de esfuerzo por riesgo}

Este anexo estima la exposición de esfuerzo de los 35 riesgos de T-16. Para cada riesgo se presentan la probabilidad de ocurrencia $p_0$, los episodios condicionales previstos $n$ y los impactos acumulados optimista (O), más probable (M) y pesimista (P), en HH. Cada escenario suma las estimaciones de los dos perfiles para un episodio y las multiplica por $n$. La estimación PERT acumulada es $\mathrm{PERT}=(O+4M+P)/6$ y el valor esperado es $\mathrm{VE}=p_0\times\mathrm{PERT}$. Son estimaciones de planificación, no frecuencias históricas ni compromisos de esfuerzo asignado; deben validarse con los responsables. La suma de VE se presenta como exposición inicial de esfuerzo conciliada con la proyección y reservas propuestas de 8.H, no como resultado de una simulación ni como distribución temporal o por perfil. Las respuestas no modifican la estimación inicial $p_0$ de esta tabla; la proyección posterior se calcula por separado en 8.F.

\begingroup
\renewcommand{\thetable}{8.B.1}
\renewcommand{\arraystretch}{0.9}
\begin{landscape}
\setlength{\textwidth}{\linewidth}
\begin{tablalafrox}{Estimación acumulada de esfuerzo y valor esperado por riesgo (HH). Fuente: Riesgos de formulario T-16}{tab:sd8-anexo-b-pert}{G{0.07} G{0.08} G{0.08} G{0.10} G{0.10} G{0.10} G{0.12} G{0.12}}{\cab{ID} & \cab{$p_0$} & \cab{Episodios $n$} & \cab{O (HH)} & \cab{M (HH)} & \cab{P (HH)} & \cab{PERT (HH)} & \cab{VE (HH)}}
R-001 & 45\% & 1 & 40,00 & 96,00 & 200,00 & 104,00 & 46,80\\
R-002 & 60\% & 1 & 72,00 & 184,00 & 360,00 & 194,67 & 116,80\\
R-003 & 50\% & 1 & 48,00 & 120,00 & 240,00 & 128,00 & 64,00\\
R-004 & 55\% & 1 & 40,00 & 104,00 & 208,00 & 110,67 & 60,87\\
R-005 & 50\% & 1 & 40,00 & 88,00 & 184,00 & 96,00 & 48,00\\
R-006 & 55\% & 1 & 40,00 & 96,00 & 184,00 & 101,33 & 55,73\\
R-007 & 35\% & 1 & 24,00 & 72,00 & 176,00 & 81,33 & 28,47\\
R-008 & 10\% & 6 & 12,00 & 30,00 & 72,00 & 34,00 & 3,40\\
R-009 & 35\% & 1 & 32,00 & 72,00 & 144,00 & 77,33 & 27,07\\
R-010 & 30\% & 1 & 24,00 & 64,00 & 128,00 & 68,00 & 20,40\\
R-011 & 35\% & 1 & 20,00 & 40,00 & 80,00 & 43,33 & 15,17\\
R-012 & 25\% & 1 & 20,00 & 56,00 & 112,00 & 59,33 & 14,83\\
R-013 & 55\% & 1 & 48,00 & 120,00 & 240,00 & 128,00 & 70,40\\
R-014 & 20\% & 1 & 40,00 & 96,00 & 192,00 & 102,67 & 20,53\\
R-015 & 30\% & 1 & 32,00 & 80,00 & 160,00 & 85,33 & 25,60\\
R-016 & 45\% & 1 & 32,00 & 80,00 & 160,00 & 85,33 & 38,40\\
R-017 & 25\% & 8 & 24,00 & 48,00 & 144,00 & 60,00 & 15,00\\
R-018 & 5\% & 2 & 12,00 & 24,00 & 64,00 & 28,67 & 1,43\\
R-019 & 30\% & 1 & 24,00 & 64,00 & 136,00 & 69,33 & 20,80\\
R-020 & 15\% & 2 & 16,00 & 40,00 & 80,00 & 42,67 & 6,40\\
R-021 & 15\% & 1 & 8,00 & 20,00 & 40,00 & 21,33 & 3,20\\
R-022 & 10\% & 1 & 8,00 & 20,00 & 40,00 & 21,33 & 2,13\\
R-023 & 10\% & 1 & 8,00 & 24,00 & 48,00 & 25,33 & 2,53\\
R-024 & 45\% & 3 & 18,00 & 36,00 & 72,00 & 39,00 & 17,55\\
R-025 & 35\% & 2 & 32,00 & 80,00 & 160,00 & 85,33 & 29,87\\
R-026 & 20\% & 1 & 32,00 & 80,00 & 160,00 & 85,33 & 17,07\\
R-027 & 40\% & 1 & 24,00 & 56,00 & 120,00 & 61,33 & 24,53\\
R-028 & 45\% & 1 & 24,00 & 64,00 & 128,00 & 68,00 & 30,60\\
R-029 & 55\% & 6 & 96,00 & 192,00 & 384,00 & 208,00 & 114,40\\
R-030 & 50\% & 1 & 48,00 & 128,00 & 280,00 & 140,00 & 70,00\\
R-031 & 45\% & 1 & 24,00 & 64,00 & 160,00 & 73,33 & 33,00\\
R-032 & 40\% & 2 & 32,00 & 96,00 & 224,00 & 106,67 & 42,67\\
R-033 & 25\% & 1 & 40,00 & 120,00 & 280,00 & 133,33 & 33,33\\
R-034 & 30\% & 2 & 32,00 & 96,00 & 240,00 & 109,33 & 32,80\\
R-035 & 40\% & 2 & 24,00 & 64,00 & 160,00 & 73,33 & 29,33\\
\multicolumn{6}{r}{Total VE} & & 1183,11\\
\end{tablalafrox}
\end{landscape}

\endgroup

\noindent La cuantificación depende de probabilidades, rangos de esfuerzo y episodios condicionales estimados; no demuestra frecuencias observadas ni incorpora dependencias entre riesgos. Los resultados no constituyen una simulación, una asignación por período o perfil, ni una garantía de esfuerzo suficiente. La suma publicada se redondea a dos decimales a partir de los valores por riesgo.



\subsection*{Hipótesis de ocurrencia y esfuerzo}
\addcontentsline{toc}{subsection}{Hipótesis de ocurrencia y esfuerzo}

Las probabilidades y HH de los riesgos R-028--R-035 son juicios explícitos de planificación basados en las condiciones citadas, no frecuencias observadas ni compromisos aceptados. $p_0$ significa al menos un episodio del evento durante la ventana indicada; $n$ estima episodios condicionados a que el evento ocurra. Los escenarios O/M/P de 8.B.1 son acumulados: se suman los dos perfiles por episodio y se multiplican por $n$ una sola vez, antes de aplicar PERT y $p_0$. El 38\% de rotación anual y las 2.600 entregas del peak no se usan como probabilidades.

Las HH representan trabajo extraordinario de LafroX por encima del alcance ordinario; no valorizan pérdidas de mercadería, reconstrucción del ERP del CLIENTE, tareas ordinarias de capacitación ni aumentos de demanda ya incluidos en el dimensionamiento. R-031 recoge sólo coordinación y recuperación de programación; la verificación y remediación de capacidad corresponde al control técnico de escalabilidad CT-14. Un incidente recibe un ID causal principal y sus HH se imputan una sola vez; los efectos en otros riesgos quedan vinculados sin repetir el esfuerzo. R-017 no duplica el mismo corte WAN como caída de ERP, SII, POS o stock central; R-034 cuantifica la dependencia del ERP del CLIENTE; CT-16 establece la corrección y verificación de la consolidación a cargo de LafroX. R-032 y R-033 cubren cambios excepcionales externos de soporte o restricciones efectivas de salida. El mantenimiento, la portabilidad y las actualizaciones ordinarias son obligaciones de LafroX; no se presupuestan otra vez como contingencia. Los controles técnicos se realizan dentro del trabajo ordinario y no se presupuestan nuevamente como contingencia.

R-017 adopta $p_0=25\%$ por la exposición inicial de Concepción con un enlace sin respaldo verificado y la falla del último camino utilizable hasta mes 56. Es una hipótesis, no una frecuencia inferida del caso ni de los cortes de Talca. Se estiman ocho episodios condicionados, escenarios acumulados 24/48/144 HH, PERT de 60 HH y VE de 15 HH. R-014 adopta $p_0=20\%$ y un impacto de nivel 4 por la consecuencia R18-16 y la jubilación dentro del horizonte; el evento incierto es la transferencia insuficiente de conocimiento, no el retiro conocido.

\begingroup
\renewcommand{\thetable}{8.B.2}
\begin{landscape}
\setlength{\textwidth}{\linewidth}
\begin{tablalafrox}{Hipótesis de ocurrencia, ventana y esfuerzo extraordinario de R-028--R-035. Fuente: condiciones citadas y estimación de planificación de LafroX.}{tab:sd8-ampliacion-hipotesis}{G{0.07} F{0.19} Y F{0.28} F{0.18}}{\cab{ID / $p_0$ / $n$} & \cab{Ventana de exposición} & \cab{Fundamento del juicio de probabilidad} & \cab{Base O/M/P y exclusiones} & \cab{Fuente de la condición}}
R-028 / 45\% / 1 & Desde levantamiento de finalidades hasta incorporación de cada ola de reparto, meses 1--20. & La objeción ya existe, pero el evento incierto es el desacuerdo sobre el uso propuesto. Se adopta 45\% para una ronda extraordinaria de coordinación previa a la ola; no se supone rechazo seguro. & JP y IMP: cada perfil 12/32/64 HH por episodio; un episodio. Sólo negociación y reprogramación extraordinarias, fuera del levantamiento y acuerdo ordinarios. & Caso, caps. 10 (restricción 10), 16.1 y 19; S3, S-13; S7, acuerdo 5.4.2. \\
R-029 / 55\% / 6 & Meses 13--56; desde primera marcha blanca y durante la operación nocturna. & El 38\% es rotación anual, no p0. Se propone 55\% de al menos un desborde de inducción durante 44 meses y, condicionado a ello, seis cohortes con atención extraordinaria; son hipótesis de exposición recurrente. & IMP y CAL: cada perfil 8/16/32 HH por cohorte; seis cohortes. Excluir las 2.816 HH ordinarias de 7.1.2 y el trabajo autónomo del preparador. & Caso, caps. 7.2, 10 (restricción 11) y 19; S7, pp. 21--22 y 32; T-14, 7.1.2; T-15, meses 13--56. \\
R-030 / 50\% / 1 & Levantamiento del mes 1 hasta integración y migraciones de ambas etapas, antes de H5/H10. & Se adopta 50\% ante ausencia documental confirmada y capacidad de acceso aún por verificar, para un bloqueo conjunto de descubrimiento; no se confunde documentación ausente con interfaz inexistente. & ARQ y DAT: cada perfil 24/64/140 HH; un episodio de descubrimiento y reproceso extraordinarios. Excluir levantamiento y adaptadores ordinarios. & Caso, cap. 5 y cap. 19; S3, S-18, S-14 y V-16; S7, 1.2.3 y 3.3.2. \\
R-031 / 45\% / 1 & Cierre de E2 en meses 19--20; septiembre del calendario supuesto de inicio de S7, sujeto a confirmar inicio contractual. & El peak y el congelamiento son condiciones conocidas. El 45\% representa un bloqueo del cierre en esa coincidencia, no la probabilidad de que septiembre ocurra. & JP y IMP: cada perfil 12/32/80 HH; un episodio. Sólo coordinación, validación y recuperación de calendario; la corrección técnica de una saturación pertenece al control técnico de escalabilidad CT-14. & Caso, caps. 10 (restricción 8), 17.5 y 19; S7, 7.3.1 y 7.3.5, pp. 37 y 58; T-18. \\
R-032 / 40\% / 2 & Meses 1--56; ciclos de soporte y compatibilidad efectivamente publicados, sin suponer fechas de fin de soporte. & Se propone 40\% de al menos una adecuación excepcional en 56 meses y dos episodios condicionados. La vigencia del soporte y su fecha de término se verifican con la información publicada por el proveedor. & ARQ y DES: cada perfil 8/24/56 HH por episodio; dos episodios. Excluir las actualizaciones anuales ordinarias 8.2.5 y mantenimiento ya incluido. & S4, plataforma declarada; S7, 8.2.5. \\
R-033 / 25\% / 1 & Meses 1--56 y preparación de la salida contractual; no se asume que la salida extraordinaria ocurrirá. & Se propone 25\% para una restricción extraordinaria de salida o sustitución en el contrato; el uso de AWS por sí solo no demuestra bloqueo. & ARQ y SRE: cada perfil 20/60/140 HH; un episodio. Sólo resolución excepcional de restricciones; excluir actualización anual del plan y salida ordinaria contratada. & S4, servicios gestionados; S7, 8.1.4 y fase 9. \\
R-034 / 30\% / 2 & Desde pruebas de emisión y primer despacho hasta mes 56; dos episodios excepcionales condicionados. & Se propone 30\% de al menos un bloqueo de nueva emisión por ERP indisponible en el horizonte; no se deriva de la respuesta del SII ni se duplica pérdida de sala como otro evento. & ARQ y SRE: cada perfil 8/24/60 HH por episodio; dos episodios. Sólo coordinación, recuperación y validación extraordinarias de la integración; no incluye reconstruir el ERP del CLIENTE. & Caso, cap. 10, restricciones 1 y 4; S4, pp. 78, 84 y 183. \\
R-035 / 40\% / 2 & Incorporación de empresas y conductores, y reemplazos durante operación hasta mes 56. & Se propone 40\% de al menos un rechazo individual en incorporación y operación; dos episodios condicionados. El hecho de trabajar para un tercero no demuestra rechazo. & IMP y CAL: cada perfil 6/16/40 HH por episodio; dos episodios. Sólo atención extraordinaria del rechazo; excluir inducción ordinaria de 7.1.3, firma de R-011 y habilitación de R-024. & Caso, caps. 10 (restricción 6), 15, RT-22.04 y 19; S7, 5.4.1 y 7.1.3. \\
\end{tablalafrox}
\end{landscape}
\endgroup

\clearpage
\section{Anexo 8.C: Respaldo de la valoración inicial}

Se aplican las escalas de 8.1.3 a los 35 eventos del T-16. $I_{HH}$ es el nivel de esfuerzo del PERT condicional acumulado del Anexo 8.B; $I_D$ es el mayor nivel de las dimensiones no HH identificadas en la justificación; $I=\max(I_{HH},I_D)$. La justificación señala la dimensión determinante y distingue una amenaza a aceptación (nivel 4) de una consecuencia grave (nivel 5). Son juicios iniciales de planificación sobre el evento declarado, sin descuento por eficacia futura ni cálculo de riesgo residual. La probabilidad se deriva de $p_0$ del Anexo 8.B; el fundamento de la estimación se explicita en 8.B.2. Las referencias contractuales corresponden a las disposiciones ya descritas en 8.3.3; no se estiman multas ni importes.

\begingroup
\renewcommand{\thetable}{8.C.1}
\renewcommand{\arraystretch}{1.10}
\begin{tablalafrox}{Fundamento del impacto inicial de los 35 riesgos. Fuente: eventos del T-16, PERT del Anexo 8.B y escalas de 8.1.3.}{tab:sd8-impactos-respaldo}{G{0.08} G{0.17} G{0.12} G{0.12} Y}{\cab{ID} & \cab{PERT HH / $I_{HH}$} & \cab{$I_D$ / I} & \cab{$q_I$} & \cab{Dimensión y fundamento}}
R-001 & 104,00 / 3 & 4 / 4 & 0,40 & Aceptación: La certificación EDI incompleta amenaza la evidencia de E2 y sus cuatro semanas finales; la carga manual no acredita ese alcance. \\
R-002 & 194,67 / 4 & 5 / 5 & 0,80 & Calidad sanitaria: Los vínculos sanitarios sin evidencia de origen pueden ser irrecuperables; el saneamiento no reconstruye una historia que no existe. \\
R-003 & 128,00 / 3 & 3 / 3 & 0,20 & Alcance: El historial incompleto exige reproceso de la carga y excepciones de reposición; el evento declarado no demuestra pérdida definitiva del alcance. \\
R-004 & 110,67 / 3 & 4 / 4 & 0,40 & Aceptación: Las diferencias sin explicación ni ajuste aprobado bloquean la aceptación de inventario antes del corte, condición de marcha blanca del art. 17.3. \\
R-005 & 96,00 / 3 & 4 / 4 & 0,40 & Aceptación: Sin saldo inicial conciliado no puede demostrarse M8 con carga aceptada; requiere resolver la evidencia de aceptación de los arts. 18.1--18.2. \\
R-006 & 101,33 / 3 & 4 / 4 & 0,40 & Aceptación: La ausencia de antecedentes verificables impide respaldar las comparaciones de OTIF o cuadratura exigidas para aceptar el período de referencia. \\
R-007 & 81,33 / 3 & 3 / 3 & 0,20 & Alcance: La deuda viva adicional obliga a extraer y conciliar nuevamente el entregable; no se declara pérdida de datos ni incumplimiento definitivo. \\
R-008 & 34,00 / 2 & 3 / 3 & 0,20 & Plazo: El cierre tributario queda pendiente de confirmación externa y requiere seguimiento y conciliación del entregable; el silencio no acredita aprobación. \\
R-009 & 77,33 / 2 & 4 / 4 & 0,40 & Aceptación sanitaria: La falta de aprobación bloquea la certificación térmica y el cierre de marcha blanca; el evento es aprobación pendiente, sin afirmar daño sanitario consumado. \\
R-010 & 68,00 / 2 & 4 / 4 & 0,40 & Plazo: La falta de equipos bloquea instalación y pruebas de una ola y amenaza su aceptación y secuencia de implantación. \\
R-011 & 43,33 / 2 & 4 / 4 & 0,40 & Alcance: Sin acuerdo no se acredita la cobertura del transportista ni sus pruebas de la ola, por lo que se amenaza la aceptación del alcance comprometido. \\
R-012 & 59,33 / 2 & 4 / 4 & 0,40 & Aceptación sanitaria: La negativa impide integrar el termógrafo y acreditar temperatura de la ruta; amenaza la aceptación sin demostrar por sí sola daño sanitario grave. \\
R-013 & 128,00 / 3 & 4 / 4 & 0,40 & Aceptación: El registro exclusivo fuera de la solución impide acreditar adopción y puede extender el cierre de la ola y su marcha blanca. \\
R-014 & 102,67 / 3 & 4 / 4 & 0,40 & Aceptación: el conocimiento debe quedar transferido y probado antes de la jubilación para acreditar R18-16; el antecedente de caída de 6 a 9 puntos de cumplimiento amenaza continuidad y aceptación. \\
R-015 & 85,33 / 3 & 4 / 4 & 0,40 & Aceptación: Usuarios sin evaluación completa no acreditan capacitación y certificación, condición acumulativa de cierre de marcha blanca del art. 17.3. \\
R-016 & 85,33 / 3 & 4 / 4 & 0,40 & Plazo: La decisión pendiente bloquea tareas de diseño, prueba o aceptación y consume holgura, amenazando hitos que requieren decisión coordinada. \\
R-017 & 60,00 / 2 & 4 / 4 & 0,40 & Servicio: La interrupción limita servicios centrales y acumula sincronización; amenaza continuidad y niveles de servicio, sin suponer eficacia del respaldo. \\
R-018 & 28,67 / 2 & 4 / 4 & 0,40 & Servicio: Superar la autonomía deja componentes sin servicio hasta restituir energía; amenaza continuidad sin imputar pérdidas ajenas al proyecto. \\
R-019 & 69,33 / 2 & 4 / 4 & 0,40 & Aceptación: Sin aprobación de residencia no puede habilitarse ni acreditarse la recuperación propuesta; requiere decisión antes de aceptar los entornos. \\
R-020 & 42,67 / 2 & 4 / 4 & 0,40 & Contrato: La autorización pendiente enfrenta congelamiento y plazo de remediación, amenazando cumplimiento; no demuestra por sí sola pérdida de continuidad contractual. \\
R-021 & 21,33 / 1 & 2 / 2 & 0,10 & Plazo: El efecto declarado se limita a repetir casos de un recorrido dentro del paquete de validación; no se establece afectación de un hito contractual. \\
R-022 & 21,33 / 1 & 2 / 2 & 0,10 & Plazo: Se evalúa la repetición de casos de prueba interrumpidos; las consecuencias personales del siniestro se atienden por el procedimiento del transportista y no son HH del proyecto. \\
R-023 & 25,33 / 2 & 2 / 2 & 0,10 & Plazo: El efecto del proyecto es reprogramar y repetir el recorrido de validación; la valoración no incluye pérdidas de vehículo o carga. \\
R-024 & 39,00 / 2 & 3 / 3 & 0,20 & Alcance: El reemplazo requiere habilitación, formación extraordinaria y verificación de registros del entregable de implantación. \\
R-025 & 85,33 / 3 & 3 / 3 & 0,20 & Contrato: El efecto estimado es evaluar y tramitar la solicitud; el desarrollo adicional sólo existe tras aprobación y ajuste de línea base. \\
R-026 & 85,33 / 3 & 4 / 4 & 0,40 & Contrato: Una exigencia aplicable obliga a adecuar y validar controles de la solución; amenaza cumplimiento y requiere decisión contractual, sin presumir una reforma o multa concreta. \\
R-027 & 61,33 / 2 & 3 / 3 & 0,20 & Contrato: La discrepancia afecta el cálculo y la liquidación de un período del tramo variable; el evento no establece interrupción ni terminación contractual. \\
R-028 & 68,00 / 2 & 4 / 4 & 0,40 & Aceptación: el desacuerdo puede impedir incorporar a los usuarios y completar pruebas de la ola; no se presupone habilitación de vigilancia excluida. \\
R-029 & 208,00 / 4 & 4 / 4 & 0,40 & Aceptación y calidad: personal no preparado amenaza la certificación y el flujo nocturno; el esfuerzo condicional de 208 HH también alcanza nivel 4. \\
R-030 & 140,00 / 3 & 4 / 4 & 0,40 & Aceptación y plazo: sin intercambio o extracción verificados se bloquean integración y migración exigidas, aunque los datos históricos existan. \\
R-031 & 73,33 / 2 & 4 / 4 & 0,40 & Contrato y aceptación: una condición incumplida impide cerrar marcha blanca y amenaza su hito; no se compensa con otras condiciones cumplidas. \\
R-032 & 106,67 / 3 & 4 / 4 & 0,40 & Seguridad y servicio: una incompatibilidad que excede mantenimiento ordinario puede impedir conservar la versión segura y el servicio aceptado. \\
R-033 & 133,33 / 3 & 4 / 4 & 0,40 & Contrato y continuidad: una restricción no resuelta impide una transición comprobable del servicio o la entrega utilizable de datos. \\
R-034 & 109,33 / 3 & 4 / 4 & 0,40 & Servicio y contrato: impedir una guía nueva bloquea despachos y amenaza la ventana crítica, aun cuando el WMS esté recuperado. \\
R-035 & 73,33 / 2 & 4 / 4 & 0,40 & Aceptación y adopción: un conductor que no usa el flujo deja evidencia y cobertura pendientes, aun con firma de empresa y certificación previa. \\
\end{tablalafrox}
\endgroup


\clearpage
\section{Anexo 8.D: Dependencias operativas y continuidad}
Las condiciones siguientes delimitan la operación, las decisiones y las pruebas habilitantes. Los responsables deben verificar la ejecución y eficacia de los controles antes de autorizar el alcance afectado. Los meses se cuentan desde el inicio contractual. Las referencias R identifican riesgos; CT-01--CT-16 identifican controles técnicos. CT sin número designa la Contraparte Técnica.

\begingroup
\renewcommand{\thetable}{8.D.1}
\begin{landscape}
\setlength{\textwidth}{\linewidth}
\renewcommand{\arraystretch}{1.08}
\begin{tablalafrox}{Dependencias operativas, decisiones y condiciones de continuidad. Fuente: requisitos de operación del Caso 02; arquitectura, procedimientos de continuidad y programa de implantación.}{tab:sd8-dependencias-operativas}{F{0.14} Y F{0.19} F{0.31}}{\cab{Ámbito} & \cab{Condición operativa} & \cab{Responsable y contraparte} & \cab{Momento, evidencia y decisión}}
Piloto INN-01 / CT-08 & Validar catálogo con Comercial/Calidad en meses 7--9. El paquete abarca 11--15; H5 corresponde a 11--12 y el piloto real a 13--15. & DAT dirige; IMP coordina; CAL verifica. & I-01A en mes 15: al menos 70\% de aciertos sobre al menos 100 avisos evaluables. Sin cumplimiento, no ampliar a otras rutas; mantener avisos informativos y alcance obligatorio. \\
Identidad / CT-15 & Distinguir manifiesto de 26 h, caché de 24 h y credenciales de 8/14 h; no autorizar por omisión ante pérdida de identidad central. & SEG responde; ARQ diseña; SRE ejecuta con TI/Operaciones. & AL-OFF-01 antes de H5/H10: corte de 24 h y relevos de 8/16 h, expiración, denegación, revocación y reconciliación. CAL bloquea el flujo fallido. \\
Emisión ERP / R-034 & El ERP del CLIENTE es el único emisor y no se modifica. La pérdida de Talca bloquea nuevas guías hasta recuperar el ERP; sólo se usan guías válidas ya emitidas. R-008 trata la confirmación externa del SII. & ARQ coordina; SRE ejecuta con TI y Operaciones. & AL-DTE-01 y AL-DR-01 antes de H5/H10: demostrar frontera de emisión y recuperación; RTO de 4 h y RPO de 15 min. No simular disponibilidad del ERP mediante otro emisor. \\
Stock / CT-16 & Conciliar UUID/sitio y excluir al escritor anterior antes de promover otra autoridad. Mantener capturas pendientes sin confirmar stock global con una copia vencida. & DAT hasta mes 21; SRE desde 22; ARQ diseña; CAL verifica. & AL-STOCK-01/AL-DR-01 antes de H5/H10: pérdida central, promoción, retorno y conciliación sin doble reserva. R-004 corresponde al inventario inicial y R-017 a la conectividad externa. \\
Pagos y servicios externos / CT-01--02 & Consultar estado antes de reintentar pagos y conciliar con clave idempotente. Ante caída de mapas, usar sólo rutas autorizadas disponibles y mantener nuevas consultas pendientes. & ARQ diseña; DES/SRE ejecutan; CAL verifica. & Pruebas de timeout antes y después del cargo, reenvío, caída y recuperación antes de H5/H10. Imputar un mismo corte WAN una sola vez a su causa principal. \\
Conocimiento de rutas / R-014 & Transferir conocimiento en meses 1--3 y verificar fecha V-14 de retiro. La prueba debe prescindir de la intervención del planificador. & IMP con Operaciones y planificador; CAL verifica. & Prueba de planificación en meses 13--15 sin el planificador. Conservar evidencia del resultado R18-16 y escalar conocimiento no transferido antes de habilitar autonomía. \\
Enlaces / R-017 & Verificar respaldo recibido en Concepción antes de aceptar la ola. Talca/Concepción: fibra y LTE, Starlink en espera caliente. Cross-docking: Starlink y LTE Entel/Claro; conmutación menor a 30 s. & SRE con TI/Operaciones; CAL verifica recepción. & Recepción antes de H3 y aceptación ampliada antes de H5. Ensayar pérdida del último camino Starlink de madrugada en Los Ángeles. La ventana local de 3 h no reduce la autonomía exigida de 24 h. \\
Energía / R-018 & Distinguir autonomía de UPS de 30 min de continuidad de Talca de 24 h; la autonomía de un componente no demuestra la de todo el servicio. & SRE con TI/Operaciones; CAL verifica. & Ensayar pérdida de suministro, transferencia, continuidad y restitución antes de aceptación. Registrar medición por componente y servicio; no acreditar autonomía sin evidencia. \\
Autorizaciones / R-016, R-019 & La aprobación previa de residencia de datos es distinta de la decisión de promover durante un incidente. El silencio no aprueba una transferencia ni una promoción. & JP y ARQ; titular y suplente autorizados del CLIENTE. & Durante promoción regional, si el titular no responde en 15 min, decide el suplente dentro del máximo de 30 min. Validar residencia antes de habilitar transferencia; conservar actas y evidencia. \\
Cobertura interior & Comprobar cobertura RF en cámaras reales con terminales, puertas y ocupación representativos; el número preliminar de AP no acredita cobertura. & ARQ/SRE con TI y Jefa de Bodega; CAL verifica. & Estudio y prueba antes de aceptar la instalación; corregir diseño y repetir si falla. La cobertura interior es una condición de instalación distinta del enlace WAN y de la asociación lote--sensor. \\
Redundancia física & Verificar pares de firewall, nodos, discos y switches, junto con autoridad de escritura, recuperación y capacidad efectiva. & ARQ diseña; SRE ejecuta; CAL verifica. & Antes de aceptar infraestructura, inyectar fallas y registrar conmutación y restitución. La redundancia declarada no reemplaza la prueba del servicio ni permite doble escritor. \\
Rotación nocturna / R-029 & El 38\% es rotación anual, no probabilidad del evento. Mantener inducción continua en meses 13--56, turno 22:00--06:00, aprendizaje máximo de dos horas y evaluación previa al trabajo autónomo. & IMP hasta mes 21; SRE desde 22; Jefa de Bodega y Operaciones aportan tutor y cobertura. & Verificar competencia de cada reemplazo y capacidad de tutoría; impedir trabajo autónomo no certificado. No imputar otra vez las HH ordinarias de capacitación ni confundir rotación con ausencia a sesiones o resistencia individual. \\
Recuperación regional & Distinguir objetivo de recuperación de 75 min (15+60) del máximo sanitario de 2 h; no consumir ese margen como disponibilidad autorizada. & SRE ejecuta; ARQ valida frontera; Calidad del CLIENTE decide aspectos sanitarios. & Pruebas AL-OFF-01, AL-DTE-01 y AL-DR-01, con 96 salidas válidas, antes de H5/H10 y semestralmente. Medir tiempos, integridad y condiciones de reanudación por procedimiento. \\
Cierre bajo peak / R-031, CT-14 & Acordar continuidad, usuarios y congelamiento antes del mes 19; mantener seguimiento diario del cierre E2 en meses 19--20. Separar coordinación del cierre de verificación técnica de capacidad. & JP con Contraparte Técnica/Operaciones; SRE y ARQ verifican capacidad. & Comprobar carga por ventana y dependencia, incluido ERP, antes de H5/H10 y preparación del peak. Preservar las cuatro semanas finales de evidencia y elevar refuerzo o cambio antes del hito. \\
Telemetría / R-028 & Acordar finalidades, acceso y controles con Operaciones, TI y sindicato antes de incorporar cada ola; sin cámaras en cabina ni control de jornada no acordado. & JP e IMP con las contrapartes funcionales; SEG verifica privacidad. & Acta previa a la ola, prueba de controles y restricción de usos no acordados. Diferenciar el desacuerdo institucional de la adopción individual; conservar evidencia de decisión. \\
\end{tablalafrox}
\end{landscape}
\endgroup

\clearpage
\section{Anexo 8.E: Controles técnicos preventivos}
Los controles CT-01--CT-16 establecen obligaciones de diseño, innovación y calidad de LafroX. Cada ficha define una solución técnica, su responsable y la evidencia necesaria para verificarla. La habilitación requiere ejecutar las pruebas y aprobar la evidencia de cumplimiento. CT sin número designa la Contraparte Técnica.

CAL verifica la evidencia y JP controla el cierre. Desde el mes 22 SRE asume las responsabilidades operativas de las familias DAT, DES e IMP conforme a la Tabla 8.1. No se promueve un flujo con defectos críticos o altos abiertos ni con pruebas habilitantes fallidas. Una función opcional que no supere sus pruebas permanece deshabilitada y se conserva el alcance obligatorio.

\begingroup
\renewcommand{\thetable}{8.E.1}
\begin{landscape}
\setlength{\textwidth}{\linewidth}
\renewcommand{\arraystretch}{1.08}
\begin{tablalafrox}{Controles técnicos: solución, responsable y verificación. Fuente: requisitos de arquitectura, seguridad, innovación y calidad de la solución.}{tab:cierre-ci}{G{0.08} F{0.25} Y F{0.34}}{\cab{Control} & \cab{Responsable} & \cab{Solución requerida} & \cab{Evidencia y condición de cierre}}
CT-01 & ARQ; DES implementa; CAL verifica. & Idempotencia persistente, consulta de estado antes de reintentar, conciliación y estado pendiente sin confirmar. & Simular timeout antes/después del cargo y reenvío de la misma clave: una sola imputación y conciliación verificable. Antes de H5/H10 en el alcance aplicable; CAL bloquea promoción del flujo fallido. \\
CT-02 & ARQ; SRE opera; CAL verifica. & Definir circuit breaker, caché de rutas autorizadas, límites y estado degradado ante caída del servicio de mapas elegido, sin dar nuevas rutas por calculadas. & Inyectar caída y recuperación del proveedor: continuidad autorizada de rutas disponibles, nuevas consultas pendientes y restitución sin pérdida. Antes de H5/H10 y tras cambio de proveedor; corregir el diseño si falla. \\
CT-03 & SEG; DES/SRE ejecutan; CAL verifica. & Aplicar mínimo privilegio, controles de autenticación y sesión del diseño; detección, revocación, conservación de evidencia y recuperación frente a credencial comprometida. & Ensayar uso indebido y revocación, incluida operación sin conexión conforme a política aprobada; informe sin hallazgos críticos/altos abiertos antes de G5. Corregir hallazgos antes de promover. \\
CT-04 & ARQ; SEG valida autenticidad; CAL verifica. & Validar origen, integridad, correlación y repetición antes de procesar EDI; poner mensajes inválidos en cuarentena sin producir pedidos o avisos. & Pruebas negativas con firma/origen alterados, mensaje repetido y correlación inválida: rechazo trazable sin efecto de negocio. Antes de H10 para EDI y regresión del alcance aplicable; corregir antes de promover. \\
CT-05 & SRE; SEG define política; CAL verifica. & Proteger almacenamiento y sesiones, aplicar enrolamiento y capacidades de bloqueo/gestión remota disponibles; revocar accesos tras pérdida y delimitar exposición. & Ensayo de terminal extraviado en línea y sin señal: datos protegidos, acceso no autorizado denegado y revocación/bloqueo al recuperar enlace. Antes de H5/H10 y tras cambio de política; no presumir borrado remoto sin conectividad. \\
CT-06 & SEG; SRE ejecuta; CAL verifica. & Aplicar límites, protección perimetral, aislamiento de portales y alertas; conservar capacidad para usuarios legítimos y continuidad de servicios críticos. & Inyectar ráfaga controlada en ambiente autorizado y medir bloqueo, latencia y disponibilidad contra umbrales de S4/T-13. Antes de H5/H10 y tras cambio perimetral; ningún ensayo ofensivo no autorizado en Producción. \\
CT-07 & SEG; SRE ejecuta; CAL verifica. & Restringir privilegios, separar funciones, preservar auditoría y respaldos y comprobar recuperación ante cambios no autorizados de una cuenta administrativa. & Ensayar cambio indebido, detección, contención y restauración con integridad demostrada; recuperar conforme RTO/RPO exigidos. Antes de H5/H10, restauraciones y pruebas de seguridad periódicas de S9; no cerrar con evidencia comprometida. \\
CT-08 & DAT dirige; IMP coordina piloto; CAL verifica. & Validar catálogo con Comercial/Calidad en 7--9, mostrar confianza y confirmar saldo en visita. Mantener alcance obligatorio; si no se valida, usar sólo aviso informativo por vida útil remanente. & Paquete 11--15: H5 en 11--12 y piloto real en 13--15. I-01A en mes 15: al menos 70\% de aciertos sobre al menos 100 avisos evaluables. Sin cumplir, no ampliar INN-01 a otras rutas; corregir y repetir. Registrar los resultados del piloto. \\
CT-09 & CAL; DES implementa; SRE aporta incidentes. & Comparar reproducción con incidente de terreno y limitar captura de batería/datos. Conservar diagnóstico convencional y regresión base si la función opcional no demuestra utilidad. & Conservar casos incidente--reproducción y medición de I-02A y consumo contra lo comprometido en S13, antes de ampliar captura; corregir captura/reproducción fallida y no acreditar el indicador sin evidencia. \\
CT-10 & SEG; DES implementa; CAL verifica. & Sanitizar por reglas fijas antes de sacar paquetes del terminal; minimizar datos y modelar amenazas. Suspender conjuntos afectados y producir muestra protegida ante filtración. & Pruebas con datos personales sembrados y campos no previstos: ningún dato personal sin sanitizar abandona el terminal; informe de revisión de amenazas y regresión. Bloquear la captura/exportación hasta aprobar, antes del primer uso y ante cada cambio de esquema. \\
CT-11 & DAT; Calidad del CLIENTE valida parámetros; CAL verifica. & Validar familias, mantener RNG-04 y no extender vencimiento por inferencia; presentar la estimación sólo en M9. Familia no validada conserva fecha impresa e historia visible. & Evidencia de I-03B por familia conforme a S13 antes de habilitar orden sugerido, con pruebas de no relajación de RNG-04. CAL bloquea familia no validada, DAT corrige y repite; la participación del CLIENTE no traslada responsabilidad de diseño. \\
CT-12 & DAT; DES implementa; CAL verifica. & Vincular lote--sensor mediante eventos de ubicación M2 y declarar tramos sin dato; impedir estimación cuando la historia esté incompleta. & Ensayar movimiento sin evento, sensor equivocado y sincronización tardía: historia incompleta visible y ninguna vida útil inferida sin fundamento. Antes de habilitar INN-03 y en regresión; corregir asociación antes de promover. \\
CT-13 & IMP; DAT/DES implementan; SEG revisa privacidad; CAL verifica. & Codiseñar hoja comprensible; aplicar umbral de agregación y supresión en cada emisión comparativa y revisión legal previa. Conservar versión de un bloque si no se entiende y bloquear bloque 4 si expone terceros. & Prueba I-05B con usuarios y pruebas de inferencia conforme S13 antes de habilitar cada bloque; informe de comprensión y supresión. Corregir y repetir antes de habilitar. \\
CT-14 & SRE; ARQ dimensiona; DES corrige; CAL verifica. & Demostrar capacidad efectiva por ventana y dependencia, incluidos ERP, colas y recursos compartidos; ajustar el dimensionamiento ofertado sin degradar requisitos. & Pruebas de carga/estrés y p95 según S4/T-13, incluyendo peak y cuello de botella ERP, antes de H5/H10 y preparación de peak. Medir capacidad, colas, errores y recuperación; bloquear promoción si incumple. La coordinación comercial del cierre sigue en R-031. \\
CT-15 & SEG; ARQ diseña; SRE ejecuta; CAL verifica. & Resolver continuidad y revocación de identidad central/local. Separar manifiesto 26 h, caché 24 h y credenciales 8/14 h; denegar permisos inválidos y conservar personal autorizado sin abrir accesos por omisión. & AL-OFF-01: 24 h de corte, relevos 8/16 h, expiración, denegación y reconciliación, antes de H5/H10 y tras cambios; CAL bloquea alcance afectado si falla. La habilitación requiere evidencia de continuidad y revocación. \\
CT-16 & DAT hasta mes 21; SRE desde 22; ARQ diseña; CAL verifica. & Mantener una sola autoridad de escritura, conciliar UUID/sitio, excluir escritor anterior y conservar capturas pendientes; no confirmar stock global con copia vencida. & AL-STOCK-01/AL-DR-01 antes de H5/H10: pérdida central, promoción, retorno y reconciliación sin doble escritor ni doble reserva. CAL bloquea flujo inconsistente y exige regresión tras corrección. R-004 y R-017 conservan causas externas distintas. \\
\end{tablalafrox}
\end{landscape}
\endgroup

\clearpage
\section{Anexo 8.F: Respuestas, esfuerzo y riesgo residual proyectado}
\label{sec:cierre-respuestas}
Las 35 fichas complementan 8.A y T-16. A es esfuerzo preventivo adicional total en la ventana, no costo de desarrollo ni pruebas ordinarias de la EDT. Se estima una revisión extraordinaria de excepción/coordinación/cierre por riesgo: 60\% del dueño y 40\% de CAL (ARQ cuando CAL es dueño), acumulados en toda la ventana. Para vigilancia ordinaria y aceptación pasiva A=0. C es esfuerzo de respuesta condicional total, repartido en esos perfiles y contenido en el PERT H0 de 8.B, no sumado de nuevo ni multiplicado otra vez por episodios n. Son hipótesis de dos perfiles, no compromisos ni rendimientos observados. JP debe contrastar estos esfuerzos con los paquetes existentes antes de aprobarlos; si una acción ya está presupuestada, se descuenta aquí y se recalculan sumas.

El plazo preventivo es el de 8.A.2 para las 28 dependencias, antes del hito/fecha allí indicado; para los siete externos se fija en cada ficha. A se ejecuta una vez como presupuesto acumulado y se programa a D-5 del primer hito afectado; si hay revisiones repetidas (p. ej. rotación), se reparte ese total, sin multiplicarlo por 44 meses ni por n. Las HH preventivas mínimas no sustituyen capacitación, mantenimiento, respaldo o controles obligatorios incluidos en el presupuesto ordinario.

La probabilidad posterior p1 es un juicio de planificación de ocurrencia del mismo evento y ventana de p0, condicionado a que se ejecuten los controles y se obtenga la evidencia indicada. No es frecuencia medida ni eficacia certificada; se mantiene la probabilidad de dependencias externas que LafroX no puede modificar. En historia sanitaria, retirada del planificador, energía, SII, siniestros y cambio normativo no se reduce p0 por describir una respuesta. Los otros p1 proponen menor ocurrencia por resolver temprano la dependencia especificada; esa hipótesis se invalida si no existe su acta o prueba. I1 conserva el impacto global inicial, incluyendo consecuencias sanitarias y contractuales: ni el bloqueo de salida ni una mejora de detección borran una pérdida histórica. E1 se obtiene de la banda de p1 y el factor de I1 con los mismos umbrales de 8.1.3.

El esfuerzo condicional remanente de atención es L=H0-C; no es una nueva estimación de consecuencias sanitarias ni de plazo. El esfuerzo residual esperado es p1 L. El tratamiento esperado es T=A+p1 C. El beneficio bruto en HH es B=p0 H0-p1 L. B/T compara reducción esperada de trabajo de atención con trabajo de tratamiento; el saldo neto B-T=(p0-p1)H0-A evita doble conteo. Así, costo total posterior=T+p1 L=A+p1 H0. Estos indicadores no valorizan alimentos, multas, seguridad de personas ni aceptación contractual; un retorno menor que uno se justifica por obligación o continuidad, sin eliminar el cálculo. Las respuestas contingentes no constituyen prevención adicional gratuita.

La base p0/O/M/P corresponde a las ventanas de exposición del Anexo 8.B. Los intervalos de programación de 8.H son hipótesis de distribución de HH, no cambian la exposición descrita en 8.B.2 ni convierten HH en días. D-5 y D-10 son días hábiles, salvo los minutos y turnos expresamente indicados. Los tiempos propuestos son máximos internos y se acortan cuando el SLA o el procedimiento de incidente exija actuar antes; no amplían compromisos de servicio ni autorizan indisponibilidad en ventanas críticas. Las fechas de respuesta no sustituyen los diez días hábiles contractuales de revisión/subsanación.

\begingroup
\renewcommand{\thetable}{8.F.1}
\begin{landscape}
\setlength{\textwidth}{\linewidth}
\renewcommand{\arraystretch}{1.08}
\begin{tablalafrox}{Fichas ejecutables de los 35 riesgos. Fuente: 8.A y estimaciones adicionales de planificación de S8.}{tab:cierre-fichas}{G{0.06} Y F{0.29} F{0.15} F{0.20}}{\cab{ID} & \cab{Acción preventiva y evidencia} & \cab{Contingencia y plazo desde disparador} & \cab{Perfiles / HH adicionales} & \cab{Momento y regla de programación}}
R-001 & Solicitar al Gerente Comercial, con TI y cada cadena, agenda EDI en mes 13, confirmada en mes 14, y certificación a D-5 del inicio de las cuatro semanas finales de evidencia de E2. Ante falta de respuesta, JP escala por CT; continuar pruebas independientes sin acreditar la cadena pendiente ni sustituir EDI por carga manual. Evidencia de salida: Agenda EDI y certificación de cada cadena, con pruebas de intercambio y registro de pendientes. & Al D-5 sin certificación, JP escala en un día hábil y conserva pendiente la cadena; sólo ampliar tras su certificación. & ARQ / CAL; A: 4,80/3,20; C: 9,60/6,40 HH. & Ventana de programación 13--20. A acumulado a D-5 del primer hito afectado; para repeticiones, repartir el total. Aplicar fechas de 8.A.2 y el plazo contingente de esta fila. \\
R-002 & Solicitar a la Jefa de Calidad, con TI y Bodega, evidencia de origen y decisión sobre lotes no reconstruibles; política V-07 al cierre del mes 2 y excepciones a D-5 de validar cada carga. Si no responde, JP escala por CT, se aísla el conjunto y queda pendiente su aceptación; no se inventa historia sanitaria. Evidencia de salida: Inventario de vínculos sanitarios por lote, evidencia de origen y resolución firmada de Calidad sobre cada excepción. & Plan adicional: DAT aísla el conjunto el mismo día; Calidad decide antes del corte. No cargar ni despachar lotes sin trazabilidad exigida; si no es reconstruible, elevar decisión de alcance contractual, sin inventar historia. & DAT / CAL; A: 7,20/4,80; C: 14,40/9,60 HH. & Ventana de programación 1--20. A acumulado a D-5 del primer hito afectado; para repeticiones, repartir el total. Aplicar fechas de 8.A.2 y el plazo contingente de esta fila. \\
R-003 & Solicitar al Jefe de Abastecimiento y a TI el historial de tres años y validación de V-09 al cierre del mes 3; corregir extractos a D-5 de validar la carga. Ante silencio, escalar por CT, separar períodos incompletos y continuar el perfilado disponible sin reducir la historia exigida. Evidencia de salida: Extracto de tres años perfilado y reconciliado con Abastecimiento/TI; acta V-09. & A D-5 sin extracto válido, separar períodos y pedir extracción complementaria; revalidar antes de cargar. & DAT / CAL; A: 3,60/2,40; C: 9,60/6,40 HH. & Ventana de programación 1--18. A acumulado a D-5 del primer hito afectado; para repeticiones, repartir el total. Aplicar fechas de 8.A.2 y el plazo contingente de esta fila. \\
R-004 & Solicitar a Bodega, con Operaciones y Finanzas, protocolo de conciliación en mes 2, validado en mes 3. Confirmar conteo a D-5 del corte y cerrar acta y ajustes antes de D. Si no hay respuesta, escalar por CT, repetir conciliación y bloquear el corte afectado hasta contar con evidencia aprobada. Evidencia de salida: Acta de conteo, causas de diferencias y ajustes autorizados por Bodega/Finanzas. & Bloquear corte con diferencia no explicada; DAT presenta conciliación y nueva prueba antes de la ventana autorizada. & DAT / CAL; A: 4,80/3,20; C: 12,00/8,00 HH. & Ventana de programación 1--20. A acumulado a D-5 del primer hito afectado; para repeticiones, repartir el total. Aplicar fechas de 8.A.2 y el plazo contingente de esta fila. \\
R-005 & Solicitar al Gerente de Operaciones y a Finanzas respaldos y conciliación V-08 de saldos por cliente y transportista; política en mes 2 y respuesta a D-5 de validar la carga M8 de cada ola. Ante silencio, escalar por CT y mantener M8 pendiente para esos saldos; no presumir saldo cero. Evidencia de salida: Saldos de envases aceptados por cliente y transportista, política V-08 y evidencia de M8. & Mantener saldos afectados pendientes; Finanzas resuelve antes de su carga, sin asumir saldo cero. & DAT / CAL; A: 3,60/2,40; C: 9,60/6,40 HH. & Ventana de programación 1--20. A acumulado a D-5 del primer hito afectado; para repeticiones, repartir el total. Aplicar fechas de 8.A.2 y el plazo contingente de esta fila. \\
R-006 & Solicitar a Comercial y Operaciones las horas históricas y a Finanzas los montos de rendición, con validación de V-10 al cierre del mes 1. Si no responden, escalar por CT y mantener pendiente la referencia de aceptación; la captura futura no reconstruye la historia anterior. Evidencia de salida: Fuentes históricas de horas y montos verificadas con Comercial, Operaciones y Finanzas, V-10. & Plan adicional: JP reúne áreas dentro de dos días hábiles desde bloqueo; en cinco días hábiles exige evidencia o eleva al Comité Ejecutivo la condición contractual. No sustituir historia por mediciones futuras ni aceptar resultado sin referencia. & IMP / CAL; A: 4,80/3,20; C: 9,60/6,40 HH. & Ventana de programación 1--20. A acumulado a D-5 del primer hito afectado; para repeticiones, repartir el total. Aplicar fechas de 8.A.2 y el plazo contingente de esta fila. \\
R-007 & Solicitar a Finanzas y TI el universo de deuda vigente, con respuesta inicial al cierre del mes 3 y conciliación a D-5 de validar la migración. Ante silencio, escalar por CT y mantener la cartera incompleta pendiente; preparar extracción complementaria sin omitir deuda obligatoria. Evidencia de salida: Universo de deuda vigente conciliado con Finanzas/TI y extracción complementaria comprobada. & No dar por migrada cartera incompleta; escalar a D-5 de cada carga y repetir conciliación. & DAT / CAL; A: 3,60/2,40; C: 9,60/6,40 HH. & Ventana de programación 1--18. A acumulado a D-5 del primer hito afectado; para repeticiones, repartir el total. Aplicar fechas de 8.A.2 y el plazo contingente de esta fila. \\
R-008 & Mantener trazabilidad de trámites pendientes, consulta y reintentos idempotentes hasta recibir confirmación del SII. Ante demora externa, conservar el estado pendiente y escalar según el procedimiento acordado; el silencio no equivale a aprobación. Evidencia de salida: Cola tributaria con correlación, consulta y reintento idempotente; confirmación externa conservada. & SRE registra pendiente al detectar demora, consulta dentro de 15 min y revisa cada 30 min; si afecta salida, avisa de inmediato a Finanzas/CT. No interpretar silencio como aprobación. & SRE / CAL; A: 0,00/0,00; C: 4,80/3,20 HH. & Ventana de programación 13--56. A acumulado a D-5 del primer hito afectado; para repeticiones, repartir el total. Aplicar fechas de 8.A.2 y el plazo contingente de esta fila. \\
R-009 & Solicitar a la Jefa de Calidad la matriz térmica versionada V-11 durante mes 1 y su aprobación al cierre del mes 2, antes de certificar. Si no responde, escalar por CT y mantener el tratamiento preventivo de lecturas fuera de ficha; no cerrar marcha blanca ni aprobar parámetros por el CLIENTE. Evidencia de salida: Matriz térmica V-11 versionada y firmada por Calidad por familia antes de certificación. & Sin parámetros aprobados, bloquear certificación del alcance afectado y escalar a D-5; no relajar bloqueo térmico. & CAL / ARQ; A: 3,60/2,40; C: 7,20/4,80 HH. & Ventana de programación 1--20. A acumulado a D-5 del primer hito afectado; para repeticiones, repartir el total. Aplicar fechas de 8.A.2 y el plazo contingente de esta fila. \\
R-010 & Solicitar al Jefe de TI, con Operaciones, calendario de equipos en mes 1, entrega confirmada a D-5 y recepción técnica antes de la instalación o prueba D. Ante silencio o faltante, escalar por CT, continuar trabajo independiente y bloquear las actividades sin equipo. Evidencia de salida: Recepción técnica por equipo y sitio, con conformidad de TI/Operaciones antes de la ola. & A D-5 sin equipos, JP decide reprogramación de trabajo independiente en un día hábil; no certificar con sustitutos no validados. & IMP / CAL; A: 2,40/1,60; C: 7,20/4,80 HH. & Ventana de programación 1--20. A acumulado a D-5 del primer hito afectado; para repeticiones, repartir el total. Aplicar fechas de 8.A.2 y el plazo contingente de esta fila. \\
R-011 & Solicitar al Gerente de Operaciones y al representante de cada transportista el acuerdo y la nómina desde mes 2, firmados a D-5 de su incorporación a la ola. Ante silencio, escalar por CT y no habilitar a esa empresa; conservar su cobertura pendiente sin retirarla del alcance. Evidencia de salida: Acuerdo operativo firmado y nómina de cada transportista antes de incorporar la ola. & A D-5 sin firma, escalar por CT el mismo día; no habilitar empresa pendiente ni reducir cobertura obligatoria. & IMP / CAL; A: 2,40/1,60; C: 4,80/3,20 HH. & Ventana de programación 2--20. A acumulado a D-5 del primer hito afectado; para repeticiones, repartir el total. Aplicar fechas de 8.A.2 y el plazo contingente de esta fila. \\
R-012 & Solicitar a Operaciones y al propietario permiso de instalación por vehículo con D-10 de la visita y obtenerlo a D-5. Ante silencio, escalar por CT, suspender la intervención y proponer otra visita o vehículo autorizado; mantener pendiente la evidencia térmica. Evidencia de salida: Autorización escrita por vehículo y visita; propietario y Operaciones confirman disponibilidad. & A D-5 sin permiso, suspender visita y acordar alternativa autorizada antes de instalar. & IMP / CAL; A: 2,40/1,60; C: 7,20/4,80 HH. & Ventana de programación 1--20. A acumulado a D-5 del primer hito afectado; para repeticiones, repartir el total. Aplicar fechas de 8.A.2 y el plazo contingente de esta fila. \\
R-013 & Solicitar a Comercial y Operaciones, con Bodega por turno, responsables y agenda de acompañamiento a D-5 del inicio de cada ola, convocando a D-10. Si no responden o persiste el registro externo, escalar por CT, focalizar acompañamiento y mantener pendiente la evidencia de adopción. Evidencia de salida: Transacciones de muestra registradas en el flujo nuevo y cobertura certificada por turno/ola. & Si persiste flujo paralelo, IMP abre refuerzo en el siguiente turno y mide diariamente; sin evidencia de adopción no cerrar ola. & IMP / CAL; A: 4,80/3,20; C: 12,00/8,00 HH. & Ventana de programación 13--56. A acumulado a D-5 del primer hito afectado; para repeticiones, repartir el total. Aplicar fechas de 8.A.2 y el plazo contingente de esta fila. \\
R-014 & Solicitar al Gerente de Operaciones y al planificador confirmación de retiro y disponibilidad V-14 al cierre del mes 1; transferir y validar reglas con el ayudante en meses 1--3. Probar las rutas y excepciones sin el planificador en la marcha blanca de meses 13--15 (S7, 7.2.1). Si falta respuesta o se retira antes, escalar por CT y reconstruir con ayudante e historia verificable, sin dar por validadas reglas pendientes. Evidencia de salida: V-14, reglas y excepciones transferidas al ayudante y prueba en 13--15 sin planificador. & Retiro confirmado no cambia p0 del fallo de transferencia. Si faltan reglas, IMP inicia reconstrucción en dos días hábiles y valida antes de cerrar el piloto. & IMP / CAL; A: 3,60/2,40; C: 9,60/6,40 HH. & Ventana de programación 1--15. A acumulado a D-5 del primer hito afectado; para repeticiones, repartir el total. Aplicar fechas de 8.A.2 y el plazo contingente de esta fila. \\
R-015 & Solicitar a Operaciones, Bodega, Comercial y TI, según perfil, nómina y liberación horaria a D-5 de cada sesión convocada a D-10. Ante silencio o ausencia, escalar por CT y reprogramar evaluación; usar sólo suplentes certificados y no acreditar preparación ni cierre de ola sin evidencia. Evidencia de salida: Lista de convocados, asistencia, evaluación y suplentes certificados por sesión. & Ausencia: avisar a jefatura el mismo día y fijar recuperación dentro de cinco días hábiles, antes de trabajo autónomo. & IMP / CAL; A: 2,40/1,60; C: 7,20/4,80 HH. & Ventana de programación 1--56. A acumulado a D-5 del primer hito afectado; para repeticiones, repartir el total. Aplicar fechas de 8.A.2 y el plazo contingente de esta fila. \\
R-016 & Solicitar a CT área decisora y suplente, con agenda al cierre del mes 1; cursar insumos a D-10 y pedir respuesta a D-5. Las aceptaciones formales conservan diez días hábiles. Para promover la región de contingencia, avisar al suplente a los 15 min sin respuesta del titular y obtener decisión dentro de 30 min (S4, 4.3.2.5); no esperar el plazo ordinario. Mantener bloqueada la acción sin autorización; el silencio no permite cambios. Evidencia de salida: Matriz de decisión con titular y suplente; recepción de insumos y resolución trazables. & D-5 sin decisión: escalar un día hábil después; promoción regional: suplente a 15 min y decisión dentro de 30 min, sin actuar por silencio. & JP / CAL; A: 2,40/1,60; C: 7,20/4,80 HH. & Ventana de programación 1--56. A acumulado a D-5 del primer hito afectado; para repeticiones, repartir el total. Aplicar fechas de 8.A.2 y el plazo contingente de esta fila. \\
R-017 & Mitigar. Verificar respaldos antes de aceptar la ola de Concepción; no considerar disponible un enlace sólo por estar ofertado. Mantener y ensayar los enlaces ya ofertados: en Talca y Concepción, fibra y LTE con Starlink en espera caliente; en cross-docking, Starlink con LTE Entel/Claro y conmutación menor a 30 s. En Los Ángeles de madrugada, probar pérdida del camino Starlink restante y continuidad local de su ventana de 3 h. Coordinar restitución y conciliación sin duplicar el mismo corte como falla de SII, POS o mapas. Referencia: S4, 4.2.5.3. Evidencia de salida: Recepción de respaldos por sitio y ensayo de pérdida del último camino, incluida Concepción y Los Ángeles. & SRE activa continuidad al corte y avisa a TI/CT en 15 min; comprobar conmutación menor a 30 s, ventana local y restitución antes de normalizar. & SRE / CAL; A: 2,40/1,60; C: 7,20/4,80 HH. & Ventana de programación 1--56. A acumulado a D-5 del primer hito afectado; para repeticiones, repartir el total. Aplicar fechas de 8.A.2 y el plazo contingente de esta fila. \\
R-018 & Verificar autonomía, alarmas y procedimiento de recuperación de los componentes de LafroX, coordinando las dependencias de energía con el CLIENTE/proveedor. Si se excede la autonomía, activar recuperación y registrar la indisponibilidad e impacto. Evidencia de salida: Autonomía y transferencia verificadas por componente; coordinación de energía con TI/Operaciones. & Al corte, SRE mide autonomía restante y alerta en 15 min; activar recuperación antes de agotarla y bloquear operación insegura. & SRE / CAL; A: 0,00/0,00; C: 4,80/3,20 HH. & Ventana de programación 13--56. A acumulado a D-5 del primer hito afectado; para repeticiones, repartir el total. Aplicar fechas de 8.A.2 y el plazo contingente de esta fila. \\
R-019 & Solicitar por CT, con TI y la autoridad del CLIENTE encauzada por Gerencia General, aprobación escrita de residencia y transferencia en mes 1, resuelta al cierre del mes 2 o antes de habilitar datos, lo que ocurra primero. Ante silencio, escalar y preparar alternativa de arquitectura; mantener deshabilitada la transferencia no aprobada. Evidencia de salida: Aprobación escrita de residencia/transferencia para regiones propuestas antes de habilitar datos. & Sin aprobación al mes 2, JP escala en un día hábil y ARQ presenta alternativa en cinco días hábiles; no transferir datos sin autorización. & ARQ / CAL; A: 2,40/1,60; C: 7,20/4,80 HH. & Ventana de programación 1--2. A acumulado a D-5 del primer hito afectado; para repeticiones, repartir el total. Aplicar fechas de 8.A.2 y el plazo contingente de esta fila. \\
R-020 & Solicitar a CT, Operaciones y TI flujo y suplentes al cierre del mes 1; ante incidente, pedir autorización escrita antes del cambio. Sin respuesta, escalar en una hora o antes si lo exige el plazo de atención, manteniendo contención autorizada y estado seguro. El silencio no permite intervenir durante congelamiento. Evidencia de salida: Flujo de cambios urgentes y suplentes aprobados; autorización escrita por intervención. & Durante incidente escalar en una hora o antes si exige el servicio; contener sin intervención prohibida y conservar congelamientos. & SEG / CAL; A: 2,40/1,60; C: 7,20/4,80 HH. & Ventana de programación 1--56. A acumulado a D-5 del primer hito afectado; para repeticiones, repartir el total. Aplicar fechas de 8.A.2 y el plazo contingente de esta fila. \\
R-021 & Detener de forma segura el recorrido afectado, registrar los casos de prueba incompletos y coordinar con el CLIENTE/transportista una nueva ventana antes de repetir la validación. Evidencia de salida: Ventana y recorrido de prueba confirmados por Operaciones/transportista antes de salir. & Al aviso de falla, suspender prueba, notificar el mismo día y acordar recorrido dentro de dos días hábiles; repetir sólo con vehículo seguro. & IMP / CAL; A: 0,00/0,00; C: 4,80/3,20 HH. & Ventana de programación 1--20. A acumulado a D-5 del primer hito afectado; para repeticiones, repartir el total. Aplicar fechas de 8.A.2 y el plazo contingente de esta fila. \\
R-022 & Priorizar la atención del accidente conforme a los procedimientos del transportista y del CLIENTE; suspender la prueba, registrar su impacto y reprogramar sólo cuando exista autorización para continuar. Evidencia de salida: Autorización de recorrido y protocolo de seguridad del CLIENTE/transportista. & Suspender de inmediato ante siniestro; coordinar el mismo día y fijar nueva ventana dentro de dos días hábiles, condicionada a autorización segura. & IMP / CAL; A: 0,00/0,00; C: 4,80/3,20 HH. & Ventana de programación 1--20. A acumulado a D-5 del primer hito afectado; para repeticiones, repartir el total. Aplicar fechas de 8.A.2 y el plazo contingente de esta fila. \\
R-023 & Coordinar con el CLIENTE y el transportista la respuesta al asalto, robo o incendio, proteger a las personas y registrar qué evidencia y pruebas se perdieron. Acordar una nueva ejecución segura antes de continuar la validación. Evidencia de salida: Custodia y seguridad del transporte corresponden al CLIENTE/transportista; evidencia de prueba preservada. & Interrumpir ante asalto, avisar el mismo día y acordar reejecución dentro de dos días hábiles; no reanudar sin habilitación segura. & IMP / CAL; A: 0,00/0,00; C: 4,80/3,20 HH. & Ventana de programación 1--20. A acumulado a D-5 del primer hito afectado; para repeticiones, repartir el total. Aplicar fechas de 8.A.2 y el plazo contingente de esta fila. \\
R-024 & Solicitar a Operaciones y al transportista nómina y suplentes a D-5 de incorporar la empresa y confirmar conductor, vehículo y viaje antes de cada salida. Si falta respuesta o llega un sustituto, escalar, verificar en andén y solicitar conductor certificado; no habilitar uso no autorizado ni acreditar cobertura incompleta. Evidencia de salida: Identidad, autorización, capacitación y vínculo persona--vehículo--viaje verificados en andén. & Sustituto no certificado: bloquear acceso al detectar y solicitar otro conductor ese mismo turno; verificar antes de salida. & IMP / CAL; A: 2,40/1,60; C: 7,20/4,80 HH. & Ventana de programación 2--56. A acumulado a D-5 del primer hito afectado; para repeticiones, repartir el total. Aplicar fechas de 8.A.2 y el plazo contingente de esta fila. \\
R-025 & Solicitar por CT a la gerencia peticionaria solicitud escrita y decisión del cambio antes de su primera actividad dependiente; registrar y remitir el mismo día hábil. Ante silencio, mantener línea base y continuar trabajo independiente; si amenaza un hito, escalar al Comité Ejecutivo sin iniciar desarrollo no aprobado. Evidencia de salida: Solicitud y análisis de impacto aprobados antes de iniciar trabajo fuera de alcance. & Registrar solicitud el mismo día; JP responde o escala en dos días hábiles y mantiene línea base hasta decisión formal. & JP / CAL; A: 2,40/1,60; C: 7,20/4,80 HH. & Ventana de programación 1--56. A acumulado a D-5 del primer hito afectado; para repeticiones, repartir el total. Aplicar fechas de 8.A.2 y el plazo contingente de esta fila. \\
R-026 & Revisar cambios normativos aplicables, determinar su efecto con las áreas competentes y tramitar ajustes de controles, retención o eliminación mediante decisión y control de cambios documentados, sin posponer una obligación legal aplicable. Evidencia de salida: Seguimiento de cambios aplicables con asesoría competente; decisión contractual documentada. & ARQ registra aviso el mismo día hábil, solicita interpretación dentro de dos y presenta impacto dentro de cinco tras recibirla; no inventar multas o reformas. & ARQ / CAL; A: 0,00/0,00; C: 9,60/6,40 HH. & Ventana de programación 1--56. A acumulado a D-5 del primer hito afectado; para repeticiones, repartir el total. Aplicar fechas de 8.A.2 y el plazo contingente de esta fila. \\
R-027 & Solicitar a Finanzas, con Comercial y CT, reglas de 5.4.3 aprobadas al cierre del mes 20 y tres liquidaciones en sombra reproducidas en meses 21--23 antes del pago variable del mes 24. Ante silencio, no activar el tramo; una disputa se escala y suspende sólo el devengo correspondiente según contrato, manteniendo el servicio. Evidencia de salida: Reglas de costo de servir aprobadas al mes 20 y tres liquidaciones en sombra reproducidas en 21--23. & Antes del mes 24 sin acuerdo, no activar tramo variable; Finanzas/CT reciben disputa en un día hábil y mantienen servicio mientras resuelven. & JP / CAL; A: 3,60/2,40; C: 7,20/4,80 HH. & Ventana de programación 20--56. A acumulado a D-5 del primer hito afectado; para repeticiones, repartir el total. Aplicar fechas de 8.A.2 y el plazo contingente de esta fila. \\
R-028 & Mitigar y escalar. JP coordina con Operaciones, TI y sindicato el acuerdo 5.4.2 antes de la ola: convocar a D-10 y resolver a D-5. Explicar finalidades, acceso y controles; mantener excluidas las cámaras en cabina y no usar GPS para jornada sin acuerdo previo. Ante desacuerdo, escalar por CT, mantener sólo el uso operacional autorizado y registrar la cobertura de la ola pendiente. Evidencia de salida: Acuerdo 5.4.2 con finalidades, accesos y controles comprendidos por Operaciones/TI/sindicato. & D-5 sin acuerdo: JP escala el mismo día y conserva sólo uso operacional autorizado; no cámaras en cabina ni GPS de jornada no acordado. & JP / CAL; A: 3,60/2,40; C: 7,20/4,80 HH. & Ventana de programación 1--20. A acumulado a D-5 del primer hito afectado; para repeticiones, repartir el total. Aplicar fechas de 8.A.2 y el plazo contingente de esta fila. \\
R-029 & Mitigar. Mantener 7.1.2 en meses 13--56, turno 22:00--06:00, tutor y evaluación antes del trabajo autónomo, con aprendizaje máximo de dos horas. Bodega y Operaciones entregan altas por turno antes de asignarlas. IMP conduce hasta mes 21 y SRE desde 22, familia IMP. Ante sobrecarga, usar personal certificado, reforzar la inducción y escalar por CT; no duplicar las HH ordinarias de 7.1.2. Evidencia de salida: Altas por turno, tutor, evaluación antes de autonomía y aprendizaje dentro de dos horas según 7.1.2. & Plan adicional: Bodega asigna cobertura certificada antes del turno; IMP/SRE revisa cada alta diariamente. Si supera tutoría, no asignar trabajo autónomo; presentar refuerzo en dos días hábiles y revalidar antes de habilitar. & IMP / CAL; A: 7,20/4,80; C: 19,20/12,80 HH. & Ventana de programación 13--56. A acumulado a D-5 del primer hito afectado; para repeticiones, repartir el total. Aplicar fechas de 8.A.2 y el plazo contingente de esta fila. \\
R-030 & Mitigar y escalar. ARQ solicita a TI acceso, diccionario y condiciones de extracción en mes 1 (1.2.3, 3.3.2 y V-16); DAT verifica la migración del WMS. Probar temprano sin modificar el ERP. Si falta una interfaz, evaluar intercambio de archivos controlado, sujeto a demostrar los requisitos; elevar por CT el bloqueo y sus impactos. La alternativa no acredita por sí sola latencia, integridad ni emisión tributaria. Evidencia de salida: Contrato de extracción e intercambio confirmado con TI; ensayo temprano sin modificar ERP y sustitución progresiva de WMS. & Plan adicional: ARQ eleva brecha al finalizar semana 2 del mes 1; TI entrega acceso o alternativa verificable antes de integración. Ensayar archivo controlado si satisface requisitos; bloquear certificación si persiste interfaz no verificable. & ARQ / CAL; A: 7,20/4,80; C: 14,40/9,60 HH. & Ventana de programación 1--18. A acumulado a D-5 del primer hito afectado; para repeticiones, repartir el total. Aplicar fechas de 8.A.2 y el plazo contingente de esta fila. \\
R-031 & Mitigar y escalar. JP acuerda con Operaciones y CT disponibilidad y continuidad antes de septiembre; cerrar habilitaciones en agosto según S7. En meses 19--20 mantener evidencia diaria de las seis condiciones de marcha blanca. Un defecto durante 1--25 de septiembre no autoriza despliegue: activar continuidad aprobada y R-020 si requiere autorización. Si persiste una condición incumplida, no declarar cierre; elevar plan de recuperación sin mover hitos por cuenta propia. Evidencia de salida: Disponibilidad y continuidad acordadas antes de septiembre; seguimiento diario de las seis condiciones de cierre. & Plan adicional: cerrar habilitaciones en mes 19 antes del congelamiento; JP revisa incidentes diariamente en 19--20. Escalar el mismo día una condición incumplida y acordar recuperación en dos días hábiles, sin intervención prohibida ni aceptación ficticia. & JP / CAL; A: 4,80/3,20; C: 9,60/6,40 HH. & Ventana de programación 19--20. A acumulado a D-5 del primer hito afectado; para repeticiones, repartir el total. Aplicar fechas de 8.A.2 y el plazo contingente de esta fila. \\
R-032 & Mitigar. ARQ mantiene inventario de versiones y revisa soporte y compatibilidad trimestralmente desde mes 1, con DES y SRE. Preparar y probar adecuación antes del límite verificado de soporte, dentro de ventanas autorizadas. Si el cambio excede la actualización ordinaria 8.2.5, registrar el esfuerzo extraordinario y elevar impactos por CT; usar reversión ensayada y soporte permitido, sin sostener una versión insegura por conveniencia. Evidencia de salida: Inventario de versiones y soporte, revisión trimestral y compatibilidad probada antes del límite verificado. & Cambio excepcional: ARQ registra aviso en un día hábil y presenta transición en cinco; programar antes de fin de soporte, conservando ventanas y pruebas. & ARQ / CAL; A: 3,60/2,40; C: 12,00/8,00 HH. & Ventana de programación 1--56. A acumulado a D-5 del primer hito afectado; para repeticiones, repartir el total. Aplicar fechas de 8.A.2 y el plazo contingente de esta fila. \\
R-033 & Mitigar. ARQ verifica con TI propiedad de cuentas, acceso a datos, claves, exportaciones e infraestructura como código en la recepción inicial y en cada actualización anual del plan 8.1.4. Probar una restauración utilizable, no sólo descargar archivos. Ante restricción, documentar brecha, mantener continuidad autorizada y escalar por CT alternativas y decisión contractual; no declarar reversibilidad completa ni cambiar de proveedor sin prueba y aprobación. Evidencia de salida: Cuentas, claves, exportación y restauración utilizable verificadas en recepción y revisión anual de salida. & Restricción efectiva: ARQ escala en un día hábil y presenta alternativas en cinco; no sustituir proveedor sin prueba/aprobación ni declarar reversibilidad por una descarga. & ARQ / CAL; A: 2,40/1,60; C: 14,40/9,60 HH. & Ventana de programación 1--56. A acumulado a D-5 del primer hito afectado; para repeticiones, repartir el total. Aplicar fechas de 8.A.2 y el plazo contingente de esta fila. \\
R-034 & Mitigar y escalar. ARQ/SRE acuerdan con TI y Finanzas restitución del ERP y ensayan AL-DTE-01/AL-DR-01 antes de H5/H10, incluida pérdida de sala. Mantener sólo salidas con guía preemitida aún válida y bloquear las que requieren nueva guía. Recuperar WMS en nube no restaura el ERP ni exceptúa el RTO requerido; registrar y escalar por CT cualquier incumplimiento, sin segundo emisor ni aprobación tributaria por silencio. Evidencia de salida: Procedimiento ERP con TI/Finanzas probado en AL-DTE-01/AL-DR-01 antes de H5/H10. & Al fallo ERP avisar de inmediato a TI/Finanzas, mantener sólo guías preemitidas válidas y bloquear nuevas salidas sin guía. Activar restitución con RTO 4 h/RPO 15 min exigidos; escalar a los 30 min sin recuperación encaminada. & ARQ / CAL; A: 3,60/2,40; C: 12,00/8,00 HH. & Ventana de programación 13--56. A acumulado a D-5 del primer hito afectado; para repeticiones, repartir el total. Aplicar fechas de 8.A.2 y el plazo contingente de esta fila. \\
R-035 & Mitigar y escalar. IMP coordina con Operaciones, transportista y conductor demostración y acompañamiento antes de su salida, dentro de 7.1.3. Confirmar disponibilidad a D-5 de la ola y aceptación práctica en andén. Ante rechazo, solicitar a la empresa conductor certificado y dispuesto, mantener acompañamiento y elevar por CT la cobertura pendiente; no atribuir a LafroX facultades laborales ni declarar completa la evidencia faltante. Evidencia de salida: Demostración práctica en andén e incorporación acordada con transportista y conductor habilitado. & Ante rechazo actuar antes de salida: IMP solicita conductor certificado/dispuesto ese turno y registra cobertura pendiente; refuerzo antes de siguiente asignación. & IMP / CAL; A: 3,60/2,40; C: 9,60/6,40 HH. & Ventana de programación 2--56. A acumulado a D-5 del primer hito afectado; para repeticiones, repartir el total. Aplicar fechas de 8.A.2 y el plazo contingente de esta fila. \\
\end{tablalafrox}
\end{landscape}
\endgroup

\begingroup
\renewcommand{\thetable}{8.F.2}
\begin{landscape}
\setlength{\textwidth}{\linewidth}
\renewcommand{\arraystretch}{1.08}
\begin{tablalafrox}{Riesgo residual proyectado, condicionado a evidencia; no resultado observado. Fuente: hipótesis por riesgo y escalas de S8.}{tab:cierre-residual}{G{0.06} F{0.13} F{0.17} F{0.16} Y}{\cab{ID} & \cab{p1 / nivel / factor} & \cab{I1 / factor; E1 y zona} & \cab{Esfuerzo residual} & \cab{Fundamento, evidencia y decisión posterior}}
R-001 & 25,00\% / 2 / 0,30 & 4 / 0,40; 0,12 (Media) & L=88,00 HH; p1L=22,00 HH. & Hipótesis p1=25,00\% por resolución temprana de la dependencia: Agenda EDI y certificación de cada cadena, con pruebas de intercambio y registro de pendientes. Revisar al verificar evidencia y en comité quincenal; si falla la hipótesis, conservar la valoración inicial y activar 8.F.1. \\
R-002 & 60,00\% / 4 / 0,70 & 5 / 0,80; 0,56 (Crítica) & L=170,67 HH; p1L=102,40 HH. & No se reduce p0: Inventario de vínculos sanitarios por lote, evidencia de origen y resolución firmada de Calidad sobre cada excepción. Plan adicional en 8.F.1; mantener bloqueo/escalamiento y revisar diariamente hasta cerrar evidencia. \\
R-003 & 30,00\% / 2 / 0,30 & 3 / 0,20; 0,06 (Baja) & L=112,00 HH; p1L=33,60 HH. & Hipótesis p1=30,00\% por resolución temprana de la dependencia: Extracto de tres años perfilado y reconciliado con Abastecimiento/TI; acta V-09. Revisar al verificar evidencia y en comité quincenal; si falla la hipótesis, conservar la valoración inicial y activar 8.F.1. \\
R-004 & 30,00\% / 2 / 0,30 & 4 / 0,40; 0,12 (Media) & L=90,67 HH; p1L=27,20 HH. & Hipótesis p1=30,00\% por resolución temprana de la dependencia: Acta de conteo, causas de diferencias y ajustes autorizados por Bodega/Finanzas. Revisar al verificar evidencia y en comité quincenal; si falla la hipótesis, conservar la valoración inicial y activar 8.F.1. \\
R-005 & 30,00\% / 2 / 0,30 & 4 / 0,40; 0,12 (Media) & L=80,00 HH; p1L=24,00 HH. & Hipótesis p1=30,00\% por resolución temprana de la dependencia: Saldos de envases aceptados por cliente y transportista, política V-08 y evidencia de M8. Revisar al verificar evidencia y en comité quincenal; si falla la hipótesis, conservar la valoración inicial y activar 8.F.1. \\
R-006 & 35,00\% / 3 / 0,50 & 4 / 0,40; 0,20 (Alta) & L=85,33 HH; p1L=29,87 HH. & Hipótesis p1=35,00\% por resolución temprana de la dependencia: Fuentes históricas de horas y montos verificadas con Comercial, Operaciones y Finanzas, V-10. Plan adicional en 8.F.1; mantener bloqueo/escalamiento y revisar diariamente hasta cerrar evidencia. \\
R-007 & 25,00\% / 2 / 0,30 & 3 / 0,20; 0,06 (Baja) & L=65,33 HH; p1L=16,33 HH. & Hipótesis p1=25,00\% por resolución temprana de la dependencia: Universo de deuda vigente conciliado con Finanzas/TI y extracción complementaria comprobada. Revisar al verificar evidencia y en comité quincenal; si falla la hipótesis, conservar la valoración inicial y activar 8.F.1. \\
R-008 & 10,00\% / 1 / 0,10 & 3 / 0,20; 0,02 (Baja) & L=26,00 HH; p1L=2,60 HH. & No se reduce p0: Cola tributaria con correlación, consulta y reintento idempotente; confirmación externa conservada. Revisar al verificar evidencia y en comité quincenal; si falla la hipótesis, conservar la valoración inicial y activar 8.F.1. \\
R-009 & 20,00\% / 2 / 0,30 & 4 / 0,40; 0,12 (Media) & L=65,33 HH; p1L=13,07 HH. & Hipótesis p1=20,00\% por resolución temprana de la dependencia: Matriz térmica V-11 versionada y firmada por Calidad por familia antes de certificación. Revisar al verificar evidencia y en comité quincenal; si falla la hipótesis, conservar la valoración inicial y activar 8.F.1. \\
R-010 & 20,00\% / 2 / 0,30 & 4 / 0,40; 0,12 (Media) & L=56,00 HH; p1L=11,20 HH. & Hipótesis p1=20,00\% por resolución temprana de la dependencia: Recepción técnica por equipo y sitio, con conformidad de TI/Operaciones antes de la ola. Revisar al verificar evidencia y en comité quincenal; si falla la hipótesis, conservar la valoración inicial y activar 8.F.1. \\
R-011 & 25,00\% / 2 / 0,30 & 4 / 0,40; 0,12 (Media) & L=35,33 HH; p1L=8,83 HH. & Hipótesis p1=25,00\% por resolución temprana de la dependencia: Acuerdo operativo firmado y nómina de cada transportista antes de incorporar la ola. Revisar al verificar evidencia y en comité quincenal; si falla la hipótesis, conservar la valoración inicial y activar 8.F.1. \\
R-012 & 15,00\% / 2 / 0,30 & 4 / 0,40; 0,12 (Media) & L=47,33 HH; p1L=7,10 HH. & Hipótesis p1=15,00\% por resolución temprana de la dependencia: Autorización escrita por vehículo y visita; propietario y Operaciones confirman disponibilidad. Revisar al verificar evidencia y en comité quincenal; si falla la hipótesis, conservar la valoración inicial y activar 8.F.1. \\
R-013 & 30,00\% / 2 / 0,30 & 4 / 0,40; 0,12 (Media) & L=108,00 HH; p1L=32,40 HH. & Hipótesis p1=30,00\% por resolución temprana de la dependencia: Transacciones de muestra registradas en el flujo nuevo y cobertura certificada por turno/ola. Revisar al verificar evidencia y en comité quincenal; si falla la hipótesis, conservar la valoración inicial y activar 8.F.1. \\
R-014 & 20,00\% / 2 / 0,30 & 4 / 0,40; 0,12 (Media) & L=86,67 HH; p1L=17,33 HH. & No se reduce p0: V-14, reglas y excepciones transferidas al ayudante y prueba en 13--15 sin planificador. Revisar al verificar evidencia y en comité quincenal; si falla la hipótesis, conservar la valoración inicial y activar 8.F.1. \\
R-015 & 20,00\% / 2 / 0,30 & 4 / 0,40; 0,12 (Media) & L=73,33 HH; p1L=14,67 HH. & Hipótesis p1=20,00\% por resolución temprana de la dependencia: Lista de convocados, asistencia, evaluación y suplentes certificados por sesión. Revisar al verificar evidencia y en comité quincenal; si falla la hipótesis, conservar la valoración inicial y activar 8.F.1. \\
R-016 & 25,00\% / 2 / 0,30 & 4 / 0,40; 0,12 (Media) & L=73,33 HH; p1L=18,33 HH. & Hipótesis p1=25,00\% por resolución temprana de la dependencia: Matriz de decisión con titular y suplente; recepción de insumos y resolución trazables. Revisar al verificar evidencia y en comité quincenal; si falla la hipótesis, conservar la valoración inicial y activar 8.F.1. \\
R-017 & 20,00\% / 2 / 0,30 & 4 / 0,40; 0,12 (Media) & L=48,00 HH; p1L=9,60 HH. & Hipótesis p1=20,00\% por resolución temprana de la dependencia: Recepción de respaldos por sitio y ensayo de pérdida del último camino, incluida Concepción y Los Ángeles. Revisar al verificar evidencia y en comité quincenal; si falla la hipótesis, conservar la valoración inicial y activar 8.F.1. \\
R-018 & 5,00\% / 1 / 0,10 & 4 / 0,40; 0,04 (Baja) & L=20,67 HH; p1L=1,03 HH. & No se reduce p0: Autonomía y transferencia verificadas por componente; coordinación de energía con TI/Operaciones. Revisar al verificar evidencia y en comité quincenal; si falla la hipótesis, conservar la valoración inicial y activar 8.F.1. \\
R-019 & 20,00\% / 2 / 0,30 & 4 / 0,40; 0,12 (Media) & L=57,33 HH; p1L=11,47 HH. & Hipótesis p1=20,00\% por resolución temprana de la dependencia: Aprobación escrita de residencia/transferencia para regiones propuestas antes de habilitar datos. Revisar al verificar evidencia y en comité quincenal; si falla la hipótesis, conservar la valoración inicial y activar 8.F.1. \\
R-020 & 10,00\% / 1 / 0,10 & 4 / 0,40; 0,04 (Baja) & L=30,67 HH; p1L=3,07 HH. & Hipótesis p1=10,00\% por resolución temprana de la dependencia: Flujo de cambios urgentes y suplentes aprobados; autorización escrita por intervención. Revisar al verificar evidencia y en comité quincenal; si falla la hipótesis, conservar la valoración inicial y activar 8.F.1. \\
R-021 & 15,00\% / 2 / 0,30 & 2 / 0,10; 0,03 (Baja) & L=13,33 HH; p1L=2,00 HH. & No se reduce p0: Ventana y recorrido de prueba confirmados por Operaciones/transportista antes de salir. Revisar al verificar evidencia y en comité quincenal; si falla la hipótesis, conservar la valoración inicial y activar 8.F.1. \\
R-022 & 10,00\% / 1 / 0,10 & 2 / 0,10; 0,01 (Baja) & L=13,33 HH; p1L=1,33 HH. & No se reduce p0: Autorización de recorrido y protocolo de seguridad del CLIENTE/transportista. Revisar al verificar evidencia y en comité quincenal; si falla la hipótesis, conservar la valoración inicial y activar 8.F.1. \\
R-023 & 10,00\% / 1 / 0,10 & 2 / 0,10; 0,01 (Baja) & L=17,33 HH; p1L=1,73 HH. & No se reduce p0: Custodia y seguridad del transporte corresponden al CLIENTE/transportista; evidencia de prueba preservada. Revisar al verificar evidencia y en comité quincenal; si falla la hipótesis, conservar la valoración inicial y activar 8.F.1. \\
R-024 & 30,00\% / 2 / 0,30 & 3 / 0,20; 0,06 (Baja) & L=27,00 HH; p1L=8,10 HH. & Hipótesis p1=30,00\% por resolución temprana de la dependencia: Identidad, autorización, capacitación y vínculo persona--vehículo--viaje verificados en andén. Revisar al verificar evidencia y en comité quincenal; si falla la hipótesis, conservar la valoración inicial y activar 8.F.1. \\
R-025 & 30,00\% / 2 / 0,30 & 3 / 0,20; 0,06 (Baja) & L=73,33 HH; p1L=22,00 HH. & Hipótesis p1=30,00\% por resolución temprana de la dependencia: Solicitud y análisis de impacto aprobados antes de iniciar trabajo fuera de alcance. Revisar al verificar evidencia y en comité quincenal; si falla la hipótesis, conservar la valoración inicial y activar 8.F.1. \\
R-026 & 20,00\% / 2 / 0,30 & 4 / 0,40; 0,12 (Media) & L=69,33 HH; p1L=13,87 HH. & No se reduce p0: Seguimiento de cambios aplicables con asesoría competente; decisión contractual documentada. Revisar al verificar evidencia y en comité quincenal; si falla la hipótesis, conservar la valoración inicial y activar 8.F.1. \\
R-027 & 25,00\% / 2 / 0,30 & 3 / 0,20; 0,06 (Baja) & L=49,33 HH; p1L=12,33 HH. & Hipótesis p1=25,00\% por resolución temprana de la dependencia: Reglas de costo de servir aprobadas al mes 20 y tres liquidaciones en sombra reproducidas en 21--23. Revisar al verificar evidencia y en comité quincenal; si falla la hipótesis, conservar la valoración inicial y activar 8.F.1. \\
R-028 & 25,00\% / 2 / 0,30 & 4 / 0,40; 0,12 (Media) & L=56,00 HH; p1L=14,00 HH. & Hipótesis p1=25,00\% por resolución temprana de la dependencia: Acuerdo 5.4.2 con finalidades, accesos y controles comprendidos por Operaciones/TI/sindicato. Revisar al verificar evidencia y en comité quincenal; si falla la hipótesis, conservar la valoración inicial y activar 8.F.1. \\
R-029 & 35,00\% / 3 / 0,50 & 4 / 0,40; 0,20 (Alta) & L=176,00 HH; p1L=61,60 HH. & Hipótesis p1=35,00\% por resolución temprana de la dependencia: Altas por turno, tutor, evaluación antes de autonomía y aprendizaje dentro de dos horas según 7.1.2. Plan adicional en 8.F.1; mantener bloqueo/escalamiento y revisar diariamente hasta cerrar evidencia. \\
R-030 & 35,00\% / 3 / 0,50 & 4 / 0,40; 0,20 (Alta) & L=116,00 HH; p1L=40,60 HH. & Hipótesis p1=35,00\% por resolución temprana de la dependencia: Contrato de extracción e intercambio confirmado con TI; ensayo temprano sin modificar ERP y sustitución progresiva de WMS. Plan adicional en 8.F.1; mantener bloqueo/escalamiento y revisar diariamente hasta cerrar evidencia. \\
R-031 & 35,00\% / 3 / 0,50 & 4 / 0,40; 0,20 (Alta) & L=57,33 HH; p1L=20,07 HH. & Hipótesis p1=35,00\% por resolución temprana de la dependencia: Disponibilidad y continuidad acordadas antes de septiembre; seguimiento diario de las seis condiciones de cierre. Plan adicional en 8.F.1; mantener bloqueo/escalamiento y revisar diariamente hasta cerrar evidencia. \\
R-032 & 25,00\% / 2 / 0,30 & 4 / 0,40; 0,12 (Media) & L=86,67 HH; p1L=21,67 HH. & Hipótesis p1=25,00\% por resolución temprana de la dependencia: Inventario de versiones y soporte, revisión trimestral y compatibilidad probada antes del límite verificado. Revisar al verificar evidencia y en comité quincenal; si falla la hipótesis, conservar la valoración inicial y activar 8.F.1. \\
R-033 & 20,00\% / 2 / 0,30 & 4 / 0,40; 0,12 (Media) & L=109,33 HH; p1L=21,87 HH. & Hipótesis p1=20,00\% por resolución temprana de la dependencia: Cuentas, claves, exportación y restauración utilizable verificadas en recepción y revisión anual de salida. Revisar al verificar evidencia y en comité quincenal; si falla la hipótesis, conservar la valoración inicial y activar 8.F.1. \\
R-034 & 25,00\% / 2 / 0,30 & 4 / 0,40; 0,12 (Media) & L=89,33 HH; p1L=22,33 HH. & Hipótesis p1=25,00\% por resolución temprana de la dependencia: Procedimiento ERP con TI/Finanzas probado en AL-DTE-01/AL-DR-01 antes de H5/H10. Revisar al verificar evidencia y en comité quincenal; si falla la hipótesis, conservar la valoración inicial y activar 8.F.1. \\
R-035 & 25,00\% / 2 / 0,30 & 4 / 0,40; 0,12 (Media) & L=57,33 HH; p1L=14,33 HH. & Hipótesis p1=25,00\% por resolución temprana de la dependencia: Demostración práctica en andén e incorporación acordada con transportista y conductor habilitado. Revisar al verificar evidencia y en comité quincenal; si falla la hipótesis, conservar la valoración inicial y activar 8.F.1. \\
\end{tablalafrox}
\end{landscape}
\endgroup

\begingroup
\renewcommand{\thetable}{8.F.3}
\begin{landscape}
\setlength{\textwidth}{\linewidth}
\renewcommand{\arraystretch}{1.08}
\begin{tablalafrox}{Costo-beneficio y sensibilidad de las 15 respuestas altas/críticas iniciales (HH). Fuente: modelo explícito de 8.F.}{tab:cierre-cb}{G{0.06} F{0.22} F{0.22} F{0.25} Y}{\cab{ID} & \cab{Entradas acumuladas} & \cab{Resultado esperado} & \cab{Sensibilidad / umbral} & \cab{Decisión y justificación}}
R-001 & A=8,00; C=16,00; H0=104,00; p0=45,00\%; p1=25,00\%. & B=24,80; T=12,00; B/T=2,07; neto=12,80 HH. & Sin mejora p1=p0: B/T=0,47; neto=-8,00 HH. Equilibrio neto: p1<=37,31\%. & Eficiencia de esfuerzo favorable bajo la hipótesis p1; mantener el tratamiento por obligación y no usar B/T como autorización de aceptación. \\
R-002 & A=12,00; C=24,00; H0=194,67; p0=60,00\%; p1=60,00\%. & B=14,40; T=26,40; B/T=0,55; neto=-12,00 HH. & Sin mejora p1=p0: B/T=0,55; neto=-12,00 HH. Equilibrio neto: p1<=53,84\%. & Prioridad obligatoria: proteger trazabilidad, condiciones de aceptación y continuidad; se mantiene aun con retorno HH menor que uno. Validar capacidad y elevar su esfuerzo antes del hito. \\
R-004 & A=8,00; C=20,00; H0=110,67; p0=55,00\%; p1=30,00\%. & B=33,67; T=14,00; B/T=2,40; neto=19,67 HH. & Sin mejora p1=p0: B/T=0,58; neto=-8,00 HH. Equilibrio neto: p1<=47,77\%. & Eficiencia de esfuerzo favorable bajo la hipótesis p1; mantener el tratamiento por obligación y no usar B/T como autorización de aceptación. \\
R-005 & A=6,00; C=16,00; H0=96,00; p0=50,00\%; p1=30,00\%. & B=24,00; T=10,80; B/T=2,22; neto=13,20 HH. & Sin mejora p1=p0: B/T=0,57; neto=-6,00 HH. Equilibrio neto: p1<=43,75\%. & Eficiencia de esfuerzo favorable bajo la hipótesis p1; mantener el tratamiento por obligación y no usar B/T como autorización de aceptación. \\
R-006 & A=8,00; C=16,00; H0=101,33; p0=55,00\%; p1=35,00\%. & B=25,87; T=13,60; B/T=1,90; neto=12,27 HH. & Sin mejora p1=p0: B/T=0,52; neto=-8,00 HH. Equilibrio neto: p1<=47,11\%. & Eficiencia de esfuerzo favorable bajo la hipótesis p1; mantener el tratamiento por obligación y no usar B/T como autorización de aceptación. \\
R-009 & A=6,00; C=12,00; H0=77,33; p0=35,00\%; p1=20,00\%. & B=14,00; T=8,40; B/T=1,67; neto=5,60 HH. & Sin mejora p1=p0: B/T=0,41; neto=-6,00 HH. Equilibrio neto: p1<=27,24\%. & Eficiencia de esfuerzo favorable bajo la hipótesis p1; mantener el tratamiento por obligación y no usar B/T como autorización de aceptación. \\
R-011 & A=4,00; C=8,00; H0=43,33; p0=35,00\%; p1=25,00\%. & B=6,33; T=6,00; B/T=1,06; neto=0,33 HH. & Sin mejora p1=p0: B/T=0,41; neto=-4,00 HH. Equilibrio neto: p1<=25,77\%. & Eficiencia de esfuerzo favorable bajo la hipótesis p1; mantener el tratamiento por obligación y no usar B/T como autorización de aceptación. \\
R-013 & A=8,00; C=20,00; H0=128,00; p0=55,00\%; p1=30,00\%. & B=38,00; T=14,00; B/T=2,71; neto=24,00 HH. & Sin mejora p1=p0: B/T=0,58; neto=-8,00 HH. Equilibrio neto: p1<=48,75\%. & Eficiencia de esfuerzo favorable bajo la hipótesis p1; mantener el tratamiento por obligación y no usar B/T como autorización de aceptación. \\
R-016 & A=4,00; C=12,00; H0=85,33; p0=45,00\%; p1=25,00\%. & B=20,07; T=7,00; B/T=2,87; neto=13,07 HH. & Sin mejora p1=p0: B/T=0,57; neto=-4,00 HH. Equilibrio neto: p1<=40,31\%. & Eficiencia de esfuerzo favorable bajo la hipótesis p1; mantener el tratamiento por obligación y no usar B/T como autorización de aceptación. \\
R-028 & A=6,00; C=12,00; H0=68,00; p0=45,00\%; p1=25,00\%. & B=16,60; T=9,00; B/T=1,84; neto=7,60 HH. & Sin mejora p1=p0: B/T=0,47; neto=-6,00 HH. Equilibrio neto: p1<=36,18\%. & Eficiencia de esfuerzo favorable bajo la hipótesis p1; mantener el tratamiento por obligación y no usar B/T como autorización de aceptación. \\
R-029 & A=12,00; C=32,00; H0=208,00; p0=55,00\%; p1=35,00\%. & B=52,80; T=23,20; B/T=2,28; neto=29,60 HH. & Sin mejora p1=p0: B/T=0,59; neto=-12,00 HH. Equilibrio neto: p1<=49,23\%. & Eficiencia de esfuerzo favorable bajo la hipótesis p1; mantener el tratamiento por obligación y no usar B/T como autorización de aceptación. \\
R-030 & A=12,00; C=24,00; H0=140,00; p0=50,00\%; p1=35,00\%. & B=29,40; T=20,40; B/T=1,44; neto=9,00 HH. & Sin mejora p1=p0: B/T=0,50; neto=-12,00 HH. Equilibrio neto: p1<=41,43\%. & Eficiencia de esfuerzo favorable bajo la hipótesis p1; mantener el tratamiento por obligación y no usar B/T como autorización de aceptación. \\
R-031 & A=8,00; C=16,00; H0=73,33; p0=45,00\%; p1=35,00\%. & B=12,93; T=13,60; B/T=0,95; neto=-0,67 HH. & Sin mejora p1=p0: B/T=0,47; neto=-8,00 HH. Equilibrio neto: p1<=34,09\%. & Prioridad obligatoria: proteger trazabilidad, condiciones de aceptación y continuidad; se mantiene aun con retorno HH menor que uno. Validar capacidad y elevar su esfuerzo antes del hito. \\
R-032 & A=6,00; C=20,00; H0=106,67; p0=40,00\%; p1=25,00\%. & B=21,00; T=11,00; B/T=1,91; neto=10,00 HH. & Sin mejora p1=p0: B/T=0,57; neto=-6,00 HH. Equilibrio neto: p1<=34,38\%. & Eficiencia de esfuerzo favorable bajo la hipótesis p1; mantener el tratamiento por obligación y no usar B/T como autorización de aceptación. \\
R-035 & A=6,00; C=16,00; H0=73,33; p0=40,00\%; p1=25,00\%. & B=15,00; T=10,00; B/T=1,50; neto=5,00 HH. & Sin mejora p1=p0: B/T=0,52; neto=-6,00 HH. Equilibrio neto: p1<=31,82\%. & Eficiencia de esfuerzo favorable bajo la hipótesis p1; mantener el tratamiento por obligación y no usar B/T como autorización de aceptación. \\
\end{tablalafrox}
\end{landscape}
\endgroup

\clearpage
\section{Anexo 8.G: Análisis de modos de falla y efectos (FMEA)}
FMEA aplicado a los 15 riesgos con E0 mayor o igual a 0,20. El modo de falla es el evento de T-16 y su efecto es la consecuencia allí declarada. Se utilizan cinco niveles definidos para este proyecto: severidad S=I global (1 corrección localizada, 2 ajuste de paquete, 3 reproceso sin afectar hitos, 4 amenaza a aceptación/servicio/hito, 5 consecuencia sanitaria grave o trazabilidad irrecuperable). O usa el nivel 1--5 de la misma banda de p0/p1, dentro de la ventana del evento, no una frecuencia diaria. D es dificultad para detectar antes del efecto: 1 comprobación automática integral bloqueante; 2 control documental o ensayo completo antes de liberar; 3 muestreo/reconciliación con cobertura parcial; 4 evidencia tardía o validación dispersa; 5 dato histórico no reconstruible o defecto oculto hasta el efecto.

NPR=S por O por D prioriza modos dentro de esta escala (1--125); no equivale a E, a probabilidad ni a HH. Se actúa por severidad 5 aunque NPR sea bajo. La propuesta obliga a tratar S=5, D mayor o igual a 4 y cualquier E alta/crítica. CAL conserva expediente de modo, causa, efecto, prueba y corrección. Las puntuaciones posteriores son proyecciones sujetas a la verificación de la evidencia de cada control. No se usa menor NPR para dispensar requisitos o aceptar un residual alto.

\begingroup
\renewcommand{\thetable}{8.G.1}
\begin{landscape}
\setlength{\textwidth}{\linewidth}
\renewcommand{\arraystretch}{1.08}
\begin{tablalafrox}{FMEA de 15 riesgos altos/críticos y evaluación posterior proyectada. Fuente: eventos T-16 y juicios de planificación de S8.}{tab:cierre-fmea}{G{0.06} F{0.12} F{0.22} F{0.23} F{0.18} Y}{\cab{ID} & \cab{FMEA inicial} & \cab{Modo, disparador y efecto} & \cab{Detección inicial y control} & \cab{FMEA posterior proyectada} & \cab{Acción, responsable y condición de habilitación}}
R-001 & S=4; O=3; D=4; NPR=48. & Cadena del alcance sin certificación antes del límite de 28 días. Dado que la habilitación EDI de cada cadena depende de su participación y certificación externa antes de las cuatro semanas finales de evidencia de E2, entonces puede que una cadena incluida en el alcance no complete su participación o certificación EDI antes de la fecha requerida, teniendo como consecuencia que no se acredite completo el intercambio electrónico exigido en E2 y deban reprogramarse pruebas y coordinación. & D0=4: la validación distribuida puede llegar después de la fecha necesaria. Control/acción: Agenda EDI y certificación de cada cadena, con pruebas de intercambio y registro de pendientes. & S1=4; O1=2; D1=2; NPR1=16. D1=2: acta o ensayo específico previo a habilitación; se invalida la mejora de detección si no existe evidencia. & Al D-5 sin certificación, JP escala en un día hábil y conserva pendiente la cadena; sólo ampliar tras su certificación. Dueño: ARQ; verifica CAL. \\
R-002 & S=5; O=4; D=5; NPR=100. & Lotes o vínculos históricos sin evidencia de origen al perfilar la carga. Dado que la migración sanitaria utiliza antecedentes del ERP, WMS y planillas y requiere vínculos históricos verificables con evidencia de origen, entonces puede que los datos recibidos contengan faltantes de lote o vínculos históricos que no puedan resolverse con esa evidencia, teniendo como consecuencia que la carga sanitaria afectada quede pendiente de conciliación y decisión documentada, y aumente el trabajo de saneamiento. & D0=5: faltante de origen o interfaz no verificable puede permanecer oculto hasta conciliar/aceptar. Control/acción: Inventario de vínculos sanitarios por lote, evidencia de origen y resolución firmada de Calidad sobre cada excepción. & S1=5; O1=4; D1=3; NPR1=60. D1=3: aun conciliando, la cobertura histórica puede ser parcial; no se reconstruye lo inexistente. & Plan adicional: DAT aísla el conjunto el mismo día; Calidad decide antes del corte. No cargar ni despachar lotes sin trazabilidad exigida; si no es reconstruible, elevar decisión de alcance contractual, sin inventar historia. Dueño: DAT; verifica CAL. \\
R-004 & S=4; O=4; D=4; NPR=64. & Diferencias sin causa ni ajuste aprobado en el acta previa al corte. Dado que la aceptación del inventario migrado requiere conciliar los saldos legados con el conteo físico y aprobar las explicaciones y ajustes antes del corte, entonces puede que se detecten diferencias sin explicación y ajuste aprobados a tiempo, teniendo como consecuencia que deban repetirse el saneamiento y la validación, y la carga afectada no pueda ser aceptada hasta resolverlas. & D0=4: la validación distribuida puede llegar después de la fecha necesaria. Control/acción: Acta de conteo, causas de diferencias y ajustes autorizados por Bodega/Finanzas. & S1=4; O1=2; D1=2; NPR1=16. D1=2: acta o ensayo específico previo a habilitación; se invalida la mejora de detección si no existe evidencia. & Bloquear corte con diferencia no explicada; DAT presenta conciliación y nueva prueba antes de la ventana autorizada. Dueño: DAT; verifica CAL. \\
R-005 & S=4; O=3; D=4; NPR=48. & No existe saldo inicial aceptado por cliente y transportista. Dado que la carga inicial de M8 requiere saldos de envases por cliente y transportista sustentados en antecedentes conciliables del cliente, entonces puede que los antecedentes entregados no permitan determinar un saldo conciliado por cliente y transportista, teniendo como consecuencia que deban repetirse la preparación y validación de la carga y no pueda demostrarse M8 con saldos iniciales aceptados. & D0=4: la validación distribuida puede llegar después de la fecha necesaria. Control/acción: Saldos de envases aceptados por cliente y transportista, política V-08 y evidencia de M8. & S1=4; O1=2; D1=2; NPR1=16. D1=2: acta o ensayo específico previo a habilitación; se invalida la mejora de detección si no existe evidencia. & Mantener saldos afectados pendientes; Finanzas resuelve antes de su carga, sin asumir saldo cero. Dueño: DAT; verifica CAL. \\
R-006 & S=4; O=4; D=5; NPR=80. & V-10 no obtiene horas o montos históricos verificables. Dado que las comparaciones históricas de aceptación dependen de horas verificables de entrega y montos esperados de rendición que la captura futura no reconstruye, entonces puede que falten horas verificables de entrega o montos esperados de rendición en los antecedentes históricos, teniendo como consecuencia que no puedan respaldarse con esos registros las comparaciones de OTIF o cuadratura y deba repetirse la validación del período de referencia. & D0=5: faltante de origen o interfaz no verificable puede permanecer oculto hasta conciliar/aceptar. Control/acción: Fuentes históricas de horas y montos verificadas con Comercial, Operaciones y Finanzas, V-10. & S1=4; O1=3; D1=3; NPR1=36. D1=3: aun conciliando, la cobertura histórica puede ser parcial; no se reconstruye lo inexistente. & Plan adicional: JP reúne áreas dentro de dos días hábiles desde bloqueo; en cinco días hábiles exige evidencia o eleva al Comité Ejecutivo la condición contractual. No sustituir historia por mediciones futuras ni aceptar resultado sin referencia. Dueño: IMP; verifica CAL. \\
R-009 & S=4; O=3; D=4; NPR=48. & Calidad no aprueba parámetros por producto al preparar certificación. Dado que la certificación de la parametrización térmica requiere una matriz de umbrales y duraciones por producto aprobada por Calidad del cliente, entonces puede que Calidad no apruebe esa matriz antes del cierre de las pruebas y la marcha blanca de los productos afectados, teniendo como consecuencia que no pueda completarse la evidencia de parametrización y aceptación del alcance térmico y deban reprogramarse validaciones. & D0=4: la validación distribuida puede llegar después de la fecha necesaria. Control/acción: Matriz térmica V-11 versionada y firmada por Calidad por familia antes de certificación. & S1=4; O1=2; D1=2; NPR1=16. D1=2: acta o ensayo específico previo a habilitación; se invalida la mejora de detección si no existe evidencia. & Sin parámetros aprobados, bloquear certificación del alcance afectado y escalar a D-5; no relajar bloqueo térmico. Dueño: CAL; verifica CAL. \\
R-011 & S=4; O=3; D=4; NPR=48. & Acuerdo operativo de un transportista pendiente antes de su ola. Dado que la incorporación de cada empresa transportista al reparto requiere que su acuerdo operativo esté firmado antes de la ola correspondiente, entonces puede que una empresa transportista no firme el acuerdo a tiempo, teniendo como consecuencia que sus conductores no puedan incorporarse ni completar las pruebas de esa ola y quede pendiente demostrar la cobertura del transportista. & D0=4: la validación distribuida puede llegar después de la fecha necesaria. Control/acción: Acuerdo operativo firmado y nómina de cada transportista antes de incorporar la ola. & S1=4; O1=2; D1=2; NPR1=16. D1=2: acta o ensayo específico previo a habilitación; se invalida la mejora de detección si no existe evidencia. & A D-5 sin firma, escalar por CT el mismo día; no habilitar empresa pendiente ni reducir cobertura obligatoria. Dueño: IMP; verifica CAL. \\
R-013 & S=4; O=4; D=4; NPR=64. & Transacciones se siguen registrando solo por el procedimiento anterior. Dado que la implantación exige que los usuarios incorporen sus transacciones al registro oficial de la solución, entonces puede que parte de los usuarios continúe registrando exclusivamente en el procedimiento anterior durante una ola, teniendo como consecuencia que no pueda acreditarse la cobertura de adopción requerida para cerrar esa ola y deba extenderse la intervención de implantación. & D0=4: la validación distribuida puede llegar después de la fecha necesaria. Control/acción: Transacciones de muestra registradas en el flujo nuevo y cobertura certificada por turno/ola. & S1=4; O1=2; D1=2; NPR1=16. D1=2: acta o ensayo específico previo a habilitación; se invalida la mejora de detección si no existe evidencia. & Si persiste flujo paralelo, IMP abre refuerzo en el siguiente turno y mide diariamente; sin evidencia de adopción no cerrar ola. Dueño: IMP; verifica CAL. \\
R-016 & S=4; O=3; D=4; NPR=48. & Decisión requerida no llega en la fecha acordada y bloquea la tarea. Dado que los diseños, pruebas y aceptaciones incluyen actividades que dependen de validaciones o decisiones de áreas del cliente, entonces puede que el área requerida no entregue una validación o decisión en la fecha acordada, teniendo como consecuencia que se detengan o repitan tareas dependientes, se consuma holgura y deban reprogramarse recursos. & D0=4: la validación distribuida puede llegar después de la fecha necesaria. Control/acción: Matriz de decisión con titular y suplente; recepción de insumos y resolución trazables. & S1=4; O1=2; D1=2; NPR1=16. D1=2: acta o ensayo específico previo a habilitación; se invalida la mejora de detección si no existe evidencia. & D-5 sin decisión: escalar un día hábil después; promoción regional: suplente a 15 min y decisión dentro de 30 min, sin actuar por silencio. Dueño: JP; verifica CAL. \\
R-028 & S=4; O=3; D=4; NPR=48. & Acuerdo de finalidades y controles pendiente a D-5 de la ola, u objeción formal al uso operacional propuesto.. Dado que existe objeción sindical a cámaras en cabina y al control de jornada por GPS, aunque la propuesta excluye cámaras en cabina y separa la telemetría operacional, entonces puede que no se acuerden o comprendan las finalidades y controles del uso propuesto antes de la ola de reparto, teniendo como consecuencia resistencia a la incorporación y pruebas o aceptaciones pendientes. & D0=4: la validación distribuida puede llegar después de la fecha necesaria. Control/acción: Acuerdo 5.4.2 con finalidades, accesos y controles comprendidos por Operaciones/TI/sindicato. & S1=4; O1=2; D1=2; NPR1=16. D1=2: acta o ensayo específico previo a habilitación; se invalida la mejora de detección si no existe evidencia. & D-5 sin acuerdo: JP escala el mismo día y conserva sólo uso operacional autorizado; no cámaras en cabina ni GPS de jornada no acordado. Dueño: JP; verifica CAL. \\
R-029 & S=4; O=4; D=4; NPR=64. & Preparador nuevo asignado sin tutor o evaluación, o aprendizaje que excede dos horas en el programa 7.1.2.. Dado que preparación nocturna tiene rotación anual del 38\% y requiere capacitación continua, entonces puede que las incorporaciones superen la capacidad de tutoría y evaluación o lleguen al flujo sin competencias verificadas, teniendo como consecuencia errores de preparación, refuerzo extraordinario y certificación de la ola pendiente. & D0=4: la validación distribuida puede llegar después de la fecha necesaria. Control/acción: Altas por turno, tutor, evaluación antes de autonomía y aprendizaje dentro de dos horas según 7.1.2. & S1=4; O1=3; D1=2; NPR1=24. D1=2: acta o ensayo específico previo a habilitación; se invalida la mejora de detección si no existe evidencia. & Plan adicional: Bodega asigna cobertura certificada antes del turno; IMP/SRE revisa cada alta diariamente. Si supera tutoría, no asignar trabajo autónomo; presentar refuerzo en dos días hábiles y revalidar antes de habilitar. Dueño: IMP; verifica CAL. \\
R-030 & S=4; O=3; D=5; NPR=60. & No hay contrato de interfaz o extracción verificado al cierre del mes 1, o el ensayo temprano contradice formato, acceso o comportamiento declarados.. Dado que el CLIENTE no tiene documentación técnica de las interfaces del ERP ni del WMS de 2013, que se reemplaza progresivamente, entonces puede que el levantamiento no permita verificar a tiempo contratos de intercambio y extracción del legado, teniendo como consecuencia reproceso de adaptadores o cargas y bloqueo de pruebas de integración o migración. & D0=5: faltante de origen o interfaz no verificable puede permanecer oculto hasta conciliar/aceptar. Control/acción: Contrato de extracción e intercambio confirmado con TI; ensayo temprano sin modificar ERP y sustitución progresiva de WMS. & S1=4; O1=3; D1=2; NPR1=24. D1=2: acta o ensayo específico previo a habilitación; se invalida la mejora de detección si no existe evidencia. & Plan adicional: ARQ eleva brecha al finalizar semana 2 del mes 1; TI entrega acceso o alternativa verificable antes de integración. Ensayar archivo controlado si satisface requisitos; bloquear certificación si persiste interfaz no verificable. Dueño: ARQ; verifica CAL. \\
R-031 & S=4; O=3; D=4; NPR=48. & Incidente alto/crítico, evidencia o certificación pendiente en las semanas finales del peak, sin ventana de intervención o validación disponible.. Dado que el peak de septiembre casi duplica entregas durante tres semanas y coincide con el cierre de la marcha blanca E2 en el calendario supuesto de S7, entonces puede que un incidente, una validación pendiente o la falta de disponibilidad de usuarios no pueda resolverse dentro de las ventanas autorizadas, teniendo como consecuencia incumplimiento de condiciones de cierre y necesidad de recuperación de la programación. & D0=4: la validación distribuida puede llegar después de la fecha necesaria. Control/acción: Disponibilidad y continuidad acordadas antes de septiembre; seguimiento diario de las seis condiciones de cierre. & S1=4; O1=3; D1=2; NPR1=24. D1=2: acta o ensayo específico previo a habilitación; se invalida la mejora de detección si no existe evidencia. & Plan adicional: cerrar habilitaciones en mes 19 antes del congelamiento; JP revisa incidentes diariamente en 19--20. Escalar el mismo día una condición incumplida y acordar recuperación en dos días hábiles, sin intervención prohibida ni aceptación ficticia. Dueño: JP; verifica CAL. \\
R-032 & S=4; O=3; D=4; NPR=48. & Aviso de soporte/compatibilidad que deja una versión necesaria fuera de la política aprobada o hace fallar la regresión de la actualización prevista.. Dado que la propuesta opera durante 56 meses con componentes como Laravel 13, PHP 8.5 y Proxmox, entonces puede que cambien su soporte o compatibilidad de forma que las actualizaciones ordinarias no basten, teniendo como consecuencia una adecuación extraordinaria y revalidación para mantener seguridad y continuidad. & D0=4: la validación distribuida puede llegar después de la fecha necesaria. Control/acción: Inventario de versiones y soporte, revisión trimestral y compatibilidad probada antes del límite verificado. & S1=4; O1=2; D1=2; NPR1=16. D1=2: acta o ensayo específico previo a habilitación; se invalida la mejora de detección si no existe evidencia. & Cambio excepcional: ARQ registra aviso en un día hábil y presenta transición en cinco; programar antes de fin de soporte, conservando ventanas y pruebas. Dueño: ARQ; verifica CAL. \\
R-035 & S=4; O=3; D=4; NPR=48. & Conductor asignado y habilitado rechaza el flujo digital o entrega sin la evidencia requerida, sin que la causa sea firma pendiente o reemplazo sin aviso.. Dado que los conductores de terceros no son trabajadores de Puelche y requieren incorporación acordada con su transportista, entonces puede que un conductor autorizado y convocado rechace utilizar el terminal o el flujo digital aun existiendo acuerdo de la empresa, teniendo como consecuencia entregas sin evidencia digital completa y cobertura de adopción de la ola pendiente. & D0=4: la validación distribuida puede llegar después de la fecha necesaria. Control/acción: Demostración práctica en andén e incorporación acordada con transportista y conductor habilitado. & S1=4; O1=2; D1=2; NPR1=16. D1=2: acta o ensayo específico previo a habilitación; se invalida la mejora de detección si no existe evidencia. & Ante rechazo actuar antes de salida: IMP solicita conductor certificado/dispuesto ese turno y registra cobertura pendiente; refuerzo antes de siguiente asignación. Dueño: IMP; verifica CAL. \\
\end{tablalafrox}
\end{landscape}
\endgroup

\clearpage
\section{Anexo 8.H: Reservas, perfiles y programación del esfuerzo}
\label{sec:cierre-reservas}
La exposición inicial de los 35 riesgos es de 1.183,11 HH, calculada con las probabilidades y los esfuerzos condicionales del Anexo 8.B. Los controles técnicos de diseño, innovación y calidad forman parte del trabajo ordinario; una diferencia de esfuerzo requiere estimación y aprobación específica.

La propuesta de contingencia específica es la suma p1 H0 de los 35 eventos: incluye atención contingente C y remanente L=H0-C. La prevención A es trabajo planificado adicional y se solicita como tal, separado de reserva. No se suman otra vez la exposición inicial, p1L y T=A+p1C a p1H0. La estimación corresponde al valor esperado; no determina percentiles ni suficiencia estadística de la reserva. Las hipótesis p1 dependen de evidencia; si no se verifica, se retiene la exposición inicial para replanificar y escalar, sin gastar una reserva proyectada como si ya estuviera aprobada.

La reserva protegida de S7, Tabla 7.4, pp. 45--46, es 3.072 HH: 1.152 en marcha blanca E1 (13--15) y 1.920 en producción/estabilización E1 (16--20). Ya integra 219.727,83 HH programadas sobre 216.655,83 HH base/cobertura. No se suma nuevamente a esos totales ni se presta a E2. El crédito inicial contra la contingencia específica es cero y sólo podrá reconocerse tras verificar la coincidencia de capacidades, perfiles, causas y períodos. JP podrá reconocer solapamiento sólo con la triple correspondencia riesgo--paquete--perfil/mes y acta de capacidad; al hacerlo, se resta exactamente una vez de la propuesta adicional, conservando la protección E1. Las fichas preventivas también se descuentan si ya están incluidas en paquetes.

La reserva de gestión propuesta es de 160 HH, sujeta a autorización del Comité Ejecutivo: cuatro ciclos extraordinarios por causas todavía no identificadas, uno por bloque 1--12, 13--20, 21--38 y 39--56. Cada ciclo reúne JP 8, ARQ 16, CAL 8 y SRE 8 HH: 4 por 40=160 HH. Son tolerancias de revisión, diagnóstico, coordinación y replanificación de un asunto emergente; no incluyen remediación de un riesgo conocido, controles técnicos ordinarios, cambios normativos ya registrados, correcciones ordinarias ni los 35 eventos. Si el asunto se identifica como un riesgo de T-16 se reclasifica, se anula su doble imputación y se usa su contingencia específica. Si excede el sobre de 40 HH, el Comité Ejecutivo decide capacidad/cambio antes de comprometer más trabajo. Sensibilidad de tolerancia: dos ciclos por bloque exigirían 320 HH, sin adoptar ese importe ni asegurar cobertura de eventos desconocidos.

La tabla mensual distribuye p1H0 uniformemente entre los meses de atención propuestos de cada ficha, 60\% al dueño y 40\% al apoyo (CAL, o ARQ si CAL es dueño). No predice fechas de ocurrencia ni cambia la ventana probabilística. Las cantidades de cálculo conservan precisión; el redondeo de impresión a 0,01 HH puede crear diferencias de centésimas al sumar celdas; los totales se calculan sobre los valores sin redondear. Prevención y gestión tienen sus propios calendarios y no están incluidas en esa tabla. La reserva distribuida es exposición esperada, no dotación comprometida: JP debe recalcular la ruta crítica y nivelar perfiles antes de solicitar su aprobación. No se convierten HH directamente a días.

\begingroup
\renewcommand{\thetable}{8.H.1}
\begin{landscape}
\setlength{\textwidth}{\linewidth}
\renewcommand{\arraystretch}{1.08}
\begin{tablalafrox}{Conciliación de esfuerzo y reservas, HH. Fuente: 8.B/8.F, S7 Tabla 7.4 y escenario explícito de gestión.}{tab:cierre-reserva}{F{0.28} G{0.14} Y}{\cab{Concepto} & \cab{HH} & \cab{Tratamiento y no duplicación}}
Exposición inicial vigente & 1.183,11 & Valoración previa al tratamiento; no sumar a la proyección posterior. \\
Trabajo preventivo adicional A & 180,00 & Trabajo planificado, 60/40 por ficha; descontar coincidencia con EDT antes de aprobar. \\
Contingencia específica propuesta p1H0 & 821,13 & Incluye C y L; condición: evidencia que sustenta p1. \\
Dentro de contingencia: p1C & 137,20 & Subconjunto; no reserva adicional. \\
Dentro de contingencia: p1L & 683,93 & Subconjunto de contingencia, no sumar otra vez. \\
Esfuerzo esperado posterior A+p1H0 & 1001,13 & Prevención más contingencia específica; sin gestión ni protección S7. \\
Gestión propuesta & 160,00 & Cuatro ciclos extraordinarios de 40 HH, con perfiles y bloques definidos. \\
Protección ya programada en S7 & 3.072,00 & Sólo E1; 1.152+1.920. Crédito demostrado inicialmente cero. \\
\end{tablalafrox}
\end{landscape}
\endgroup

\begingroup
\renewcommand{\thetable}{8.H.2}
\begin{landscape}
\setlength{\textwidth}{\linewidth}
\renewcommand{\arraystretch}{1.08}
\begin{tablalafrox}{Trabajo preventivo adicional propuesto por mes de primera atención y perfil (HH). Fuente: fichas 8.F.1; ejecutar antes del primer hito dependiente.}{tab:cierre-preventiva}{G{0.07} G{0.10} G{0.10} G{0.10} G{0.10} G{0.10} G{0.10} G{0.10}}{\cab{Mes} & \cab{JP} & \cab{ARQ} & \cab{DAT} & \cab{CAL} & \cab{SRE} & \cab{IMP} & \cab{SEG}}
1 & 8,40 & 18,00 & 22,80 & 48,40 & 2,40 & 15,60 & 2,40 \\
2 & 0,00 & 0,00 & 0,00 & 5,60 & 0,00 & 8,40 & 0,00 \\
13 & 0,00 & 8,40 & 0,00 & 13,60 & 0,00 & 12,00 & 0,00 \\
19 & 4,80 & 0,00 & 0,00 & 3,20 & 0,00 & 0,00 & 0,00 \\
20 & 3,60 & 0,00 & 0,00 & 2,40 & 0,00 & 0,00 & 0,00 \\
\end{tablalafrox}
\end{landscape}
\endgroup

\begingroup
\renewcommand{\thetable}{8.H.3}
\begin{landscape}
\setlength{\textwidth}{\linewidth}
\renewcommand{\arraystretch}{1.08}
\begin{tablalafrox}{Distribución mensual propuesta de contingencia específica (HH). Fuente: p1H0 y ventanas de 8.F.}{tab:cierre-mensual}{G{0.06} G{0.09} G{0.09} G{0.09} G{0.09} G{0.09} G{0.09} G{0.09} G{0.10}}{\cab{Mes} & \cab{JP} & \cab{ARQ} & \cab{DAT} & \cab{CAL} & \cab{SRE} & \cab{IMP} & \cab{SEG} & \cab{Total}}
1 & 1,01 & 6,86 & 7,32 & 12,49 & 0,13 & 2,98 & 0,05 & 30,83 \\
2 & 1,01 & 6,86 & 7,32 & 12,93 & 0,13 & 3,65 & 0,05 & 31,95 \\
3 & 1,01 & 2,70 & 7,32 & 10,16 & 0,13 & 3,65 & 0,05 & 25,02 \\
4 & 1,01 & 2,70 & 7,32 & 10,16 & 0,13 & 3,65 & 0,05 & 25,02 \\
5 & 1,01 & 2,70 & 7,32 & 10,16 & 0,13 & 3,65 & 0,05 & 25,02 \\
6 & 1,01 & 2,70 & 7,32 & 10,16 & 0,13 & 3,65 & 0,05 & 25,02 \\
7 & 1,01 & 2,70 & 7,32 & 10,16 & 0,13 & 3,65 & 0,05 & 25,02 \\
8 & 1,01 & 2,70 & 7,32 & 10,16 & 0,13 & 3,65 & 0,05 & 25,02 \\
9 & 1,01 & 2,70 & 7,32 & 10,16 & 0,13 & 3,65 & 0,05 & 25,02 \\
10 & 1,01 & 2,70 & 7,32 & 10,16 & 0,13 & 3,65 & 0,05 & 25,02 \\
11 & 1,01 & 2,70 & 7,32 & 10,16 & 0,13 & 3,65 & 0,05 & 25,02 \\
12 & 1,01 & 2,70 & 7,32 & 10,16 & 0,13 & 3,65 & 0,05 & 25,02 \\
13 & 1,01 & 5,02 & 7,32 & 12,76 & 0,19 & 5,17 & 0,05 & 31,52 \\
14 & 1,01 & 5,02 & 7,32 & 12,76 & 0,19 & 5,17 & 0,05 & 31,52 \\
15 & 1,01 & 5,02 & 7,32 & 12,76 & 0,19 & 5,17 & 0,05 & 31,52 \\
16 & 1,01 & 5,02 & 7,32 & 12,22 & 0,19 & 4,34 & 0,05 & 30,16 \\
17 & 1,01 & 5,02 & 7,32 & 12,22 & 0,19 & 4,34 & 0,05 & 30,16 \\
18 & 1,01 & 5,02 & 7,32 & 12,22 & 0,19 & 4,34 & 0,05 & 30,16 \\
19 & 8,71 & 3,39 & 5,36 & 14,96 & 0,19 & 4,34 & 0,05 & 37,00 \\
20 & 8,96 & 3,39 & 5,36 & 15,12 & 0,19 & 4,34 & 0,05 & 37,42 \\
21 & 0,75 & 1,13 & 0,00 & 2,76 & 0,19 & 2,03 & 0,05 & 6,91 \\
22 & 0,75 & 1,13 & 0,00 & 2,76 & 0,19 & 2,03 & 0,05 & 6,91 \\
23 & 0,75 & 1,13 & 0,00 & 2,76 & 0,19 & 2,03 & 0,05 & 6,91 \\
24 & 0,75 & 1,13 & 0,00 & 2,76 & 0,19 & 2,03 & 0,05 & 6,91 \\
25 & 0,75 & 1,13 & 0,00 & 2,76 & 0,19 & 2,03 & 0,05 & 6,91 \\
26 & 0,75 & 1,13 & 0,00 & 2,76 & 0,19 & 2,03 & 0,05 & 6,91 \\
27 & 0,75 & 1,13 & 0,00 & 2,76 & 0,19 & 2,03 & 0,05 & 6,91 \\
28 & 0,75 & 1,13 & 0,00 & 2,76 & 0,19 & 2,03 & 0,05 & 6,91 \\
29 & 0,75 & 1,13 & 0,00 & 2,76 & 0,19 & 2,03 & 0,05 & 6,91 \\
30 & 0,75 & 1,13 & 0,00 & 2,76 & 0,19 & 2,03 & 0,05 & 6,91 \\
31 & 0,75 & 1,13 & 0,00 & 2,76 & 0,19 & 2,03 & 0,05 & 6,91 \\
32 & 0,75 & 1,13 & 0,00 & 2,76 & 0,19 & 2,03 & 0,05 & 6,91 \\
33 & 0,75 & 1,13 & 0,00 & 2,76 & 0,19 & 2,03 & 0,05 & 6,91 \\
34 & 0,75 & 1,13 & 0,00 & 2,76 & 0,19 & 2,03 & 0,05 & 6,91 \\
35 & 0,75 & 1,13 & 0,00 & 2,76 & 0,19 & 2,03 & 0,05 & 6,91 \\
36 & 0,75 & 1,13 & 0,00 & 2,76 & 0,19 & 2,03 & 0,05 & 6,91 \\
37 & 0,75 & 1,13 & 0,00 & 2,76 & 0,19 & 2,03 & 0,05 & 6,91 \\
38 & 0,75 & 1,13 & 0,00 & 2,76 & 0,19 & 2,03 & 0,05 & 6,91 \\
39 & 0,75 & 1,13 & 0,00 & 2,76 & 0,19 & 2,03 & 0,05 & 6,91 \\
40 & 0,75 & 1,13 & 0,00 & 2,76 & 0,19 & 2,03 & 0,05 & 6,91 \\
41 & 0,75 & 1,13 & 0,00 & 2,76 & 0,19 & 2,03 & 0,05 & 6,91 \\
42 & 0,75 & 1,13 & 0,00 & 2,76 & 0,19 & 2,03 & 0,05 & 6,91 \\
43 & 0,75 & 1,13 & 0,00 & 2,76 & 0,19 & 2,03 & 0,05 & 6,91 \\
44 & 0,75 & 1,13 & 0,00 & 2,76 & 0,19 & 2,03 & 0,05 & 6,91 \\
45 & 0,75 & 1,13 & 0,00 & 2,76 & 0,19 & 2,03 & 0,05 & 6,91 \\
46 & 0,75 & 1,13 & 0,00 & 2,76 & 0,19 & 2,03 & 0,05 & 6,91 \\
47 & 0,75 & 1,13 & 0,00 & 2,76 & 0,19 & 2,03 & 0,05 & 6,91 \\
48 & 0,75 & 1,13 & 0,00 & 2,76 & 0,19 & 2,03 & 0,05 & 6,91 \\
49 & 0,75 & 1,13 & 0,00 & 2,76 & 0,19 & 2,03 & 0,05 & 6,91 \\
50 & 0,75 & 1,13 & 0,00 & 2,76 & 0,19 & 2,03 & 0,05 & 6,91 \\
51 & 0,75 & 1,13 & 0,00 & 2,76 & 0,19 & 2,03 & 0,05 & 6,91 \\
52 & 0,75 & 1,13 & 0,00 & 2,76 & 0,19 & 2,03 & 0,05 & 6,91 \\
53 & 0,75 & 1,13 & 0,00 & 2,76 & 0,19 & 2,03 & 0,05 & 6,91 \\
54 & 0,75 & 1,13 & 0,00 & 2,76 & 0,19 & 2,03 & 0,05 & 6,91 \\
55 & 0,75 & 1,13 & 0,00 & 2,76 & 0,19 & 2,03 & 0,05 & 6,91 \\
56 & 0,75 & 1,13 & 0,00 & 2,76 & 0,19 & 2,03 & 0,05 & 6,91 \\
\end{tablalafrox}
\end{landscape}
\endgroup
\paragraph{Programación y autorización.} JP asigna las fichas preventivas y la contingencia a los paquetes de migración/datos, implantación, integración, soporte, gobierno o salida que atienden cada evento. Antes de autorizar capacidad adicional debe verificar paquete, perfil, mes y capacidad existente, descontando toda coincidencia. La prevención se programa en el mes de primera atención y antes del hito dependiente; la contingencia se nivela con la distribución mensual propuesta. Las 3.072 HH protegidas mantienen su uso exclusivo en E1. Los ciclos de gestión se sitúan en los meses 1, 13, 21 y 39 como capacidad de decisión; sólo se consumen ante un asunto emergente autorizado. Las familias DAT/IMP conservan su clasificación de esfuerzo y SRE asume responsabilidad desde mes 22, sin duplicar cargas por relevo.

La programación está sujeta a nivelación y autorización de la línea base. JP verifica precedencias, holguras y ruta crítica, preservando congelamientos y las cuatro semanas finales de evidencia. Si la capacidad no permite cumplir, eleva refuerzo o cambio antes del hito, sin desplazarlo unilateralmente. La valoración económica utiliza la misma diferencia aprobada por mes y perfil, multiplicada por la tarifa neta validada. Gestión y protección no se registran como costo consumido por el solo hecho de reservarlas. JP presenta la conciliación al arranque, la actualiza en cada comité quincenal y antes de H5/H10; el Comité Ejecutivo autoriza los cambios y CAL verifica la ausencia de doble imputación. Una propuesta de reserva no acredita financiamiento aprobado.


\section*{Referencias}
\addcontentsline{toc}{section}{Referencias}
\begingroup
\small\setlength{\parindent}{-0.5in}\setlength{\leftskip}{0.5in}\setlength{\parskip}{3pt}
Distribuidora Puelche S.A. (2026a). \emph{Bases Administrativas de Licitación N.º TFEP-01/2026: Contratación de Solución Integral de Software y Servicios de Operación} (Formulario T-16).


LafroX SpA. (2026a). \emph{Propuesta técnica: Subdocumento 4, Arquitectura lógica, arquitectura física y data center}.

LafroX SpA. (2026b). \emph{Propuesta técnica: Subdocumento 7, Plan de trabajo}, y formularios T-14 y T-15.

LafroX SpA. (2026d). \emph{Propuesta técnica: Subdocumento 13, Innovaciones}.

LafroX SpA. (2026c). \emph{Propuesta técnica: Subdocumento 9, Plan de calidad}, y formularios T-13 y T-17.

Distribuidora Puelche S.A. (2026b). \emph{Bases Técnicas del Caso 02: Logística} (caps. 5--7, 10, 14--19).


\endgroup

\end{document}
